% Lesson 4 — Listening: Texts about taking part in communal work — Anglais 1ère A — Cours signé OnBuch+
% Programme officiel MINESEC. Compiler avec : tectonic ENA04-listening-texts-about-taking-part-in-communal-work.tex
\newcommand{\DOCMATIERE}{Anglais}
\newcommand{\DOCNIVEAU}{1ère A}
\newcommand{\DOCMODULE}{Chapter 1 : Family and Social Life}
\newcommand{\DOCLECON}{ENA04}
\newcommand{\DOCTITRE}{Lesson 4 — Listening: Texts about taking part in communal work}
% OnBuch+ — préambule « cours signés » ENRICHI (compilé avec tectonic / XeLaTeX)
% Système visuel « Pop » d'OnBuch+ : fond crème (jamais blanc), bordures encre
% épaisses, ombres « dures » décalées, titres Archivo Black en capitales.
% Version MATHÉMATIQUES : graphes (pgfplots/tikz), boîtes pédagogiques
% (définition, propriété, méthode, attention, le savais-tu, exercice résolu,
% exercice, TP, bilan), page de garde et signature OnBuch+ ancrée en bas.
%
% Macros attendues dans main.tex AVANT \input{preamble} :
%   \DOCMATIERE  \DOCNIVEAU  \DOCMODULE  \DOCLECON (numéro)  \DOCTITRE
\documentclass[11pt]{article}

\usepackage{fontspec}
\usepackage[english,french]{babel} % français = langue principale ; l'anglais via \eng{..} / textebox
\usepackage[a4paper,margin=2cm,top=3.4cm,bottom=2.8cm,headheight=1.4cm,footskip=1.3cm]{geometry}
\usepackage{amsmath,amssymb,amsfonts}
\usepackage{mathtools} % \xmapsto, \coloneqq... souvent utilisés par les modèles
\usepackage{cancel} % simplifications barrées (bilans de forces)
% \abs{z} (module, valeur absolue) et \norm{u} (norme), utilisés par les modèles
\providecommand{\abs}{}\let\abs\relax\DeclarePairedDelimiter{\abs}{\lvert}{\rvert}
\providecommand{\norm}{}\let\norm\relax\DeclarePairedDelimiter{\norm}{\lVert}{\rVert}
\usepackage[table]{xcolor} % [table] : fournit aussi \rowcolors (alternance de lignes)
\usepackage{pagecolor}
\usepackage{tikz}
\usetikzlibrary{arrows.meta,positioning,calc,shapes.geometric,shapes.misc,shapes.arrows,decorations.pathmorphing,decorations.pathreplacing,decorations.markings,patterns,shadows,mindmap,trees,fit,backgrounds,matrix,intersections,angles,quotes}
\usepackage{pgfplots}
\usepackage[version=4]{mhchem} % équations nucléaires \ce{^{235}_{92}U}
\usepackage[european,siunitx]{circuitikz} % schémas électriques
\pgfplotsset{compat=1.18}
\usepgfplotslibrary{fillbetween,groupplots} % aires entre courbes (calcul intégral)
\usepackage{enumitem}
\usepackage{fancyhdr}
% [most] charge toutes les bibliothèques tcolorbox (listings, minted, raster,
% documentation...) et gonfle inutilement la mémoire TeX sur un document long
% (~300 boîtes) : on ne charge que ce qui est réellement utilisé.
\usepackage[skins,breakable]{tcolorbox}
\usepackage{graphicx}
\usepackage{colortbl}
\usepackage{booktabs}
\usepackage{tabularx}
\usepackage{diagbox} % case d'en-tête diagonale des tableaux à double entrée
% Colonnes extensibles centrée / gauche / droite (C, L, R), souvent utilisées par les modèles
\newcolumntype{C}{>{\centering\arraybackslash}X}
\newcolumntype{L}{>{\raggedright\arraybackslash}X}
\newcolumntype{R}{>{\raggedleft\arraybackslash}X}
\usepackage{array}
\usepackage{multicol}
\usepackage{multirow}
\usepackage{siunitx}
% parse-numbers=false : accepte aussi \SI{2,5 \times 10^{-3}}{...}
\sisetup{locale=FR,output-decimal-marker={,},group-minimum-digits=4,parse-numbers=false}
\DeclareSIUnit\ohm{\ensuremath{\symup{\Omega}}} % U+2126 absent de la police
\DeclareSIUnit\division{div} % réglages d'oscilloscope (ms/div, V/div)
\DeclareSIUnit\tr{tr} % tours (vitesse de rotation, ex. \SI{1500}{\tr\per\minute})
\DeclareSIUnit\molar{M} % concentration molaire (ex. \SI{3}{\molar}, \SI{1}{\micro\molar})
\DeclareSIUnit\an{an} % année (le modèle confond parfois avec \year, un compteur TeX) — ex. \SI{5}{\kilo\meter\per\an}
\DeclareSIUnit\FCFA{FCFA} % franc CFA (ex. \SI{500}{\FCFA\per\kilo\gram})
\DeclareSIUnit\degreeBrix{\textdegree Brix} % teneur en sucre d'un sirop/jus (réfractométrie)
\DeclareSIUnit\month{mois} % le modèle utilise parfois \month, un compteur TeX, comme unité de durée
\usepackage{qrcode}
% \degree : commande fréquemment utilisée par les modèles pour un angle ou une
% température en dehors de \SI{}{} (non fournie par les paquets chargés).
\providecommand{\degree}{^{\circ}}
% \qed (fin de démonstration), utilisé par les modèles sans amsthm
\providecommand{\qed}{\hfill$\square$}
% \barre : double barre des valeurs interdites dans un tableau de variations (tkz-tab)
\providecommand{\barre}{\ensuremath{\|}}
% environnement proof (amsthm non chargé) utilisé par les modèles
\ifdefined\proof\else\newenvironment{proof}[1][Démonstration]{\par\noindent\textit{#1.}\ }{\qed\par}\fi
\usepackage{lastpage}
\usepackage{hyperref}
\hypersetup{hidelinks}

% ============================================================
% Palette « Pop » OnBuch+ — mêmes hex que la classe P (plus_ui.dart)
% ============================================================
\definecolor{poporange}{HTML}{F59321}
\definecolor{popdark}{HTML}{E07A0C}
\definecolor{popcream}{HTML}{FFF6E8}
\definecolor{popink}{HTML}{1C1714}
\definecolor{poppurple}{HTML}{7A5AE0}
\definecolor{popgreen}{HTML}{2E9E6A}
\definecolor{poppink}{HTML}{E0457B}
\definecolor{popblue}{HTML}{2F7DE1}
\definecolor{popgold}{HTML}{C98A12}
\definecolor{popmuted}{HTML}{8A7F76}
\definecolor{popline}{HTML}{EADFD2}
\definecolor{popgray}{HTML}{8A7F76}
\colorlet{popgrey}{popgray}
% Alias tolérants (le modèle nomme parfois les couleurs différemment)
\colorlet{popwhite}{white}
\colorlet{popblack}{popink}
\colorlet{popred}{poppink}
\colorlet{popyellow}{popgold}
% Teintes claires (fonds de tableaux/encadrés) — le modèle nomme parfois
% les couleurs avec un suffixe L (ex. popblueL) sans les avoir définies.
\colorlet{poporangeL}{poporange!20}
\colorlet{popdarkL}{popdark!20}
\colorlet{poppurpleL}{poppurple!20}
\colorlet{popgreenL}{popgreen!20}
\colorlet{poppinkL}{poppink!20}
\colorlet{popblueL}{popblue!20}
\colorlet{popgoldL}{popgold!20}
\colorlet{popredL}{popred!20}
\colorlet{popgrayL}{popgray!20}
% Teintes claires pour fonds de tableaux / zones
\colorlet{popblueL}{popblue!10!white}
\colorlet{poporangeL}{poporange!14!white}
\colorlet{popgreenL}{popgreen!12!white}
\colorlet{poppurpleL}{poppurple!12!white}
\colorlet{poppinkL}{poppink!12!white}
\colorlet{popgoldL}{popgold!14!white}

\pagecolor{popcream}

% ============================================================
% Typographie
% ============================================================
\defaultfontfeatures{Ligatures=TeX}
\setmainfont{PlusJakartaSans-Medium.ttf}[Path=../../../fonts/,BoldFont=PlusJakartaSans-Bold.ttf]
% Symboles, API et emojis absents de la police principale : repli sur DejaVu Sans (caractères actifs)
\newfontface\symfont{DejaVuSans.ttf}[Path=../../../fonts/]
\makeatletter
\newcommand\symactive[1]{\catcode#1=13 \begingroup\lccode`\~=#1 \lowercase{\endgroup\def~}{{\symfont\char#1}}}
\newcommand\symblank[1]{\catcode#1=13 \begingroup\lccode`\~=#1 \lowercase{\endgroup\def~}{}}
\newcommand\symhyph[1]{\catcode#1=13 \begingroup\lccode`\~=#1 \lowercase{\endgroup\def~}{\mbox{-}}}
\newcount\symc
\newcommand\symrange[2]{\symc=#1 \@whilenum\symc<#2 \do{\expandafter\symactive\expandafter{\the\symc}\advance\symc by 1 }}
\newcommand\symblankrange[2]{\symc=#1 \@whilenum\symc<#2 \do{\expandafter\symblank\expandafter{\the\symc}\advance\symc by 1 }}
\makeatother
\symrange{"0250}{"02FF}\symactive{"2713}\symactive{"2714}\symactive{"2717}\symactive{"2718}\symactive{"2610}\symactive{"2611}\symactive{"2612}
\symhyph{"2011}\symhyph{"2E1A}\symblankrange{"1F300}{"1FAFF}\symblank{"FE0F}
% Maths en sans-serif (Fira Math) avec chiffres et lettres droites de Plus
% Jakarta, pour que les formules se fondent dans le texte.
\usepackage[math-style=ISO,bold-style=ISO]{unicode-math}
\setmathfont{FiraMath-Regular.otf}[Path=../../../fonts/]
\setmathfont{PlusJakartaSans-Medium.ttf}[Path=../../../fonts/,range={up/num,up/{Latin,latin},\mathup}]
\usepackage{icomma} % virgule décimale sans espace en mode math
% Caractères mathématiques Unicode absents des polices de texte (Plus Jakarta,
% Archivo Black) : rendus par leur équivalent mathématique, en texte comme en
% formule — sinon ils s'affichent en carré vide (ex. « Arithmétique dans ℤ »).
\usepackage{newunicodechar}
\newunicodechar{ℝ}{\ensuremath{\mathbb{R}}}
\newunicodechar{ℤ}{\ensuremath{\mathbb{Z}}}
\newunicodechar{ℂ}{\ensuremath{\mathbb{C}}}
\newunicodechar{ℕ}{\ensuremath{\mathbb{N}}}
\newunicodechar{ℚ}{\ensuremath{\mathbb{Q}}}
\newunicodechar{𝔻}{\ensuremath{\mathbb{D}}}
\newunicodechar{α}{\ensuremath{\alpha}}
\newunicodechar{β}{\ensuremath{\beta}}
\newunicodechar{γ}{\ensuremath{\gamma}}
\newunicodechar{δ}{\ensuremath{\delta}}
\newunicodechar{ε}{\ensuremath{\varepsilon}}
\newunicodechar{θ}{\ensuremath{\theta}}
\newunicodechar{λ}{\ensuremath{\lambda}}
\newunicodechar{μ}{\ensuremath{\mu}}
\newunicodechar{σ}{\ensuremath{\sigma}}
\newunicodechar{τ}{\ensuremath{\tau}}
\newunicodechar{φ}{\ensuremath{\varphi}}
\newunicodechar{ψ}{\ensuremath{\psi}}
\newunicodechar{ω}{\ensuremath{\omega}}
\newunicodechar{Σ}{\ensuremath{\Sigma}}
% majuscules produites par \MakeUppercase dans les titres (α-aminés -> Α)
\newunicodechar{Α}{\ensuremath{\alpha}}
\newunicodechar{Β}{\ensuremath{\beta}}
\newunicodechar{Γ}{\ensuremath{\Gamma}}
\newunicodechar{Δ}{\ensuremath{\Delta}}
\newunicodechar{Ε}{\ensuremath{\varepsilon}}
\newunicodechar{Θ}{\ensuremath{\Theta}}
\newunicodechar{Λ}{\ensuremath{\Lambda}}
\newunicodechar{Μ}{\ensuremath{\mu}}
\newunicodechar{Τ}{\ensuremath{\tau}}
\newunicodechar{Φ}{\ensuremath{\Phi}}
\newunicodechar{Ψ}{\ensuremath{\Psi}}
\newunicodechar{Ω}{\ensuremath{\Omega}}
\newunicodechar{↦}{\ensuremath{\mapsto}}
\newunicodechar{∈}{\ensuremath{\in}}
\newunicodechar{∉}{\ensuremath{\notin}}
\newunicodechar{∧}{\ensuremath{\wedge}}
\newunicodechar{⇔}{\ensuremath{\Leftrightarrow}}
\newunicodechar{⟺}{\ensuremath{\Longleftrightarrow}}
\newunicodechar{⇒}{\ensuremath{\Rightarrow}}
\newunicodechar{⟹}{\ensuremath{\Longrightarrow}}
\newunicodechar{∘}{\ensuremath{\circ}}
\newunicodechar{∪}{\ensuremath{\cup}}
\newunicodechar{∩}{\ensuremath{\cap}}
\newunicodechar{≡}{\ensuremath{\equiv}}
\newunicodechar{⊂}{\ensuremath{\subset}}
\newunicodechar{⊥}{\ensuremath{\perp}}
\newunicodechar{∃}{\ensuremath{\exists}}
\newunicodechar{∀}{\ensuremath{\forall}}
\newunicodechar{∛}{\ensuremath{\sqrt[3]{\,}}}
\newunicodechar{∜}{\ensuremath{\sqrt[4]{\,}}}
\newunicodechar{⃗}{\ensuremath{{}^{\to}}}
\newunicodechar{ⁿ}{\ifmmode^{n}\else\textsuperscript{\ensuremath{n}}\fi}
\newunicodechar{ˣ}{\ifmmode^{x}\else\textsuperscript{\ensuremath{x}}\fi}
\newunicodechar{ᵇ}{\ifmmode^{b}\else\textsuperscript{\ensuremath{b}}\fi}
\newunicodechar{ʳ}{\ifmmode^{r}\else\textsuperscript{\ensuremath{r}}\fi}
\newunicodechar{ᵘ}{\ifmmode^{u}\else\textsuperscript{\ensuremath{u}}\fi}
\newunicodechar{ʸ}{\ifmmode^{y}\else\textsuperscript{\ensuremath{y}}\fi}
\newunicodechar{ᵃ}{\ifmmode^{a}\else\textsuperscript{\ensuremath{a}}\fi}
\newunicodechar{ᵉ}{\ifmmode^{e}\else\textsuperscript{\ensuremath{e}}\fi}
\newunicodechar{ᵏ}{\ifmmode^{k}\else\textsuperscript{\ensuremath{k}}\fi}
\newunicodechar{ᵖ}{\ifmmode^{p}\else\textsuperscript{\ensuremath{p}}\fi}
\newunicodechar{ᴺ}{\ifmmode^{N}\else\textsuperscript{\ensuremath{N}}\fi}
\newunicodechar{ᶜ}{\ifmmode^{c}\else\textsuperscript{\ensuremath{c}}\fi}
\newunicodechar{ʰ}{\ifmmode^{h}\else\textsuperscript{\ensuremath{h}}\fi}
\newunicodechar{ᵠ}{\ifmmode^{q}\else\textsuperscript{\ensuremath{q}}\fi}
\newunicodechar{ₙ}{\ifmmode_{n}\else\textsubscript{\ensuremath{n}}\fi}
\newunicodechar{ₐ}{\ifmmode_{a}\else\textsubscript{\ensuremath{a}}\fi}
\newunicodechar{ᵢ}{\ifmmode_{i}\else\textsubscript{\ensuremath{i}}\fi}
\newunicodechar{ᵣ}{\ifmmode_{r}\else\textsubscript{\ensuremath{r}}\fi}
\newunicodechar{ₖ}{\ifmmode_{k}\else\textsubscript{\ensuremath{k}}\fi}
\newunicodechar{ᵦ}{\ifmmode_{\beta}\else\textsubscript{\ensuremath{\beta}}\fi}
\newunicodechar{ₓ}{\ifmmode_{x}\else\textsubscript{\ensuremath{x}}\fi}
\newfontfamily\popdisplay{ArchivoBlack-Regular.ttf}[Path=../../../fonts/]
\newfontfamily\poplogo{SpaceGrotesk-Bold.ttf}[Path=../../../fonts/]
\newfontfamily\popsemi{PlusJakartaSans-SemiBold.ttf}[Path=../../../fonts/]
\newfontfamily\popxbold{PlusJakartaSans-ExtraBold.ttf}[Path=../../../fonts/]
\newfontfamily\popxboldls{PlusJakartaSans-ExtraBold.ttf}[Path=../../../fonts/,LetterSpace=3.0]

\setlist[enumerate]{leftmargin=1.7em,itemsep=5pt,topsep=4pt}
\setlist[itemize]{leftmargin=1.7em,itemsep=4pt,topsep=4pt,label={\color{poporange}\textbullet}}
\setlength{\parindent}{0pt}
\setlength{\parskip}{5pt}
\renewcommand{\arraystretch}{1.25}

% pgfplots : style Pop par défaut
\pgfplotsset{
  every axis/.append style={axis line style={popink,line width=1pt},tick style={popink},
    grid style={popline,line width=0.5pt},label style={font=\small},tick label style={font=\footnotesize}},
  popcurve/.style={poporange,line width=1.8pt,no markers,smooth},
}
% Même style disponible dans un \draw TikZ simple (hors pgfplots)
\tikzset{popcurve/.style={poporange,line width=1.8pt}}
\tikzset{popgrid/.style={popline,very thin}}
\colorlet{popcurve}{poporange} % le nom du style est parfois utilisé comme couleur

% ============================================================
% Pill / squiggle
% ============================================================
\newcommand{\poppill}[3]{%
  \tikz[baseline=-0.55ex]\node[draw=popink,line width=1.2pt,rounded corners=2.6pt,%
    fill=#2,inner xsep=6.5pt,inner ysep=3pt]{%
    \popxboldls\fontsize{7}{7}\selectfont\color{#3}\MakeUppercase{#1}};%
}
\newcommand{\popsquiggle}[1]{%
  \tikz[baseline=-0.4ex]{
    \draw[#1,line width=2.6pt,line cap=round]
      plot [smooth] coordinates {(0,0.09) (0.34,0.01) (0.68,0.21) (1.02,0.01) (1.36,0.10)};
  }%
}
% Mot-clé surligné (marqueur orange sous le texte)
\newcommand{\cle}[1]{{\popxbold\color{popdark}#1}}

% ============================================================
% En-tête / pied
% ============================================================
\pagestyle{fancy}
\fancyhf{}
\renewcommand{\headrulewidth}{0pt}
\renewcommand{\footrulewidth}{0pt}
\lhead{\raisebox{-0.15ex}{\poplogo\fontsize{13}{13}\selectfont\color{popink}OnBuch\color{poporange}+}}
\rhead{\poppill{\DOCMATIERE\ \textbullet\ \DOCNIVEAU\ \textbullet\ Leçon \DOCLECON}{white}{popink}}
\lfoot{\popxbold\fontsize{7.6}{7.6}\selectfont\color{popmuted}\MakeUppercase{Cours signé OnBuch+}\ \textbullet\ {\popsemi\DOCTITRE}}
\rfoot{\tikz[baseline=-0.6ex]\node[rounded corners=8.5pt,fill=popink,minimum height=17pt,minimum width=17pt,inner xsep=5pt,inner sep=0]{\popxbold\fontsize{8}{8}\selectfont\color{popcream}\thepage\,/\,\pageref*{LastPage}};}

% ============================================================
% Titres
% ============================================================
\newcommand{\chaptitre}[2]{%
  \begin{tcolorbox}[enhanced,colback=poporange,colframe=popink,boxrule=2.6pt,arc=11pt,
      drop shadow=popink,left=15pt,right=15pt,top=12pt,bottom=11pt,before skip=3pt,after skip=18pt]
    \poppill{#2}{popink}{popcream}\par\vspace{9pt}
    {\raggedright\hyphenpenalty=10000\exhyphenpenalty=10000\popdisplay\fontsize{22}{24}\selectfont\color{popcream}\MakeUppercase{#1}\par}
  \end{tcolorbox}
}
% Section numérotée : pastille orange avec le numéro + titre Archivo
\newcounter{popsec}
\newcommand{\coursec}[1]{%
  \stepcounter{popsec}\setcounter{popsub}{0}%
  \par\vspace{16pt}\noindent
  \begin{minipage}{\linewidth}
  \tikz[baseline=-0.7ex]\node[draw=popink,line width=1.6pt,fill=poporange,rounded corners=4pt,minimum size=22pt,inner sep=0]{\popdisplay\fontsize{12}{12}\selectfont\color{popcream}\thepopsec};%
  \hspace{8pt}{\popdisplay\fontsize{15}{17}\selectfont\color{popink}\MakeUppercase{#1}}\par
  \vspace{4pt}\noindent{\color{poporange}\rule{36pt}{3.6pt}}
  \end{minipage}\par\nopagebreak\vspace{8pt}
}
\newcounter{popsub}
\newcommand{\courssub}[1]{%
  \stepcounter{popsub}%
  \par\vspace{9pt}\noindent{\popxbold\fontsize{11.5}{13}\selectfont\color{popink}{\color{poporange}\thepopsec.\thepopsub}\hspace{6pt}#1}\par\nopagebreak\vspace{4pt}
}

% ============================================================
% Boîtes pédagogiques — même recette : fond blanc, bordure épaisse colorée,
% ombre dure encre, badge plein. Titre optionnel [#1].
% ============================================================
\tcbset{popbase/.style={breakable,enhanced,colback=white,boxrule=2.3pt,arc=9pt,
  shadow={3pt}{-3pt}{0pt}{popink},left=14pt,right=14pt,top=11pt,bottom=11pt,before skip=11pt,after skip=15pt,
  fontupper=\popsemi\fontsize{10.6}{14.5}\selectfont\color{popink},
  fontlower=\popsemi\fontsize{10.6}{14.5}\selectfont\color{popink}}}
% Coche dessinée (le glyphe ✓ n'existe pas dans les polices du document)
\newcommand{\popcheck}[1]{\tikz[baseline=-0.5ex]\draw[#1,line width=1.6pt,line cap=round,line join=round](-0.13,0.0)--(-0.04,-0.09)--(0.13,0.1);}
\providecommand{\coche}{\popcheck{popgreen}}
\providecommand{\resultat}[1]{\par\smallskip\noindent\fcolorbox{popgreen}{popgreenL}{\parbox{\dimexpr\linewidth-2\fboxsep-2\fboxrule\relax}{\textbf{Résultat.} #1}}\par\smallskip}
\newcommand{\popboxhead}[3]{\poppill{#1}{#2}{white}\ifx\relax#3\relax\else\hspace{6pt}{\popxbold\fontsize{10.6}{12}\selectfont\color{popink}#3}\fi\par\vspace{7pt}}

\newtcolorbox{aretenir}[1][]{popbase,colframe=popgreen,before upper={\popboxhead{À retenir}{popgreen}{#1}}}
% Texte en anglais (document de lecture, dialogue, extrait) : cadre
% d'étude. Un titre optionnel (source, auteur) s'affiche dans l'en-tête.
\newtcolorbox{textebox}[1][]{popbase,colframe=popblue,before upper={\popboxhead{Text}{popblue}{#1}\begin{otherlanguage}{english}},after upper={\end{otherlanguage}}}
% Marques ✓ / ✗ (forme correcte / fautive) souvent utilisées par les modèles
\providecommand{\cross}{\textcolor{poppink}{\ensuremath{\times}}}
\providecommand{\xmark}{\cross}\providecommand{\crossmark}{\cross}\providecommand{\cmark}{\ensuremath{\checkmark}}
% Ligne de réponse / justification à compléter (inventée par les modèles)
\providecommand{\justification}[1]{\par\noindent\textit{Justification~:} \dotfill\par}
% \textipa (package tipa non chargé) : la police ne couvre pas l'API -> simple passage
\providecommand{\textipa}[1]{#1}
% Soulignement (ulem non chargé) : \uline, \ul
\providecommand{\uline}[1]{\underline{#1}}\providecommand{\ul}[1]{\underline{#1}}
% \glossaire{\item ...} inventée par les modèles : liste de glossaire
\providecommand{\glossaire}[1]{\begin{itemize}#1\end{itemize}}
% \alert (beamer) inventée par les modèles : texte mis en évidence
\providecommand{\alert}[1]{\textbf{\textcolor{poppink}{#1}}}
% variantes ASCII d'ordinaux écrites en macro
\providecommand{\iere}{\textsuperscript{re}}\providecommand{\ieme}{\textsuperscript{e}}\providecommand{\ere}{\textsuperscript{re}}\providecommand{\eme}{\textsuperscript{e}}\providecommand{\er}{\textsuperscript{er}}\providecommand{\ie}{\textsuperscript{e}}
% Ordinaux français écrits en macro par les modèles : 1\ière, 3\ième
\newcommand{\ière}{\textsuperscript{re}}\newcommand{\ième}{\textsuperscript{e}}
% \exemple (inventée par les modèles) : étiquette « Exemple : »
\providecommand{\exemple}{\textbf{Exemple~:}\ }
% \square utilisable aussi hors mode math (cases à cocher)
\AtBeginDocument{\let\squareMath\square\renewcommand{\square}{\ensuremath{\squareMath}}}
% Mot ou phrase en anglais à l'intérieur d'un paragraphe français
\newcommand{\eng}[1]{\ifmmode\text{\textit{#1}}\else\begingroup\selectlanguage{english}\itshape #1\endgroup\fi}
\let\esp\eng % alias toléré
% flèches d'intonation inventées par les modèles
\providecommand{\textsearrow}{}\renewcommand{\textsearrow}{\ensuremath{\searrow}}\providecommand{\textnearrow}{}\renewcommand{\textnearrow}{\ensuremath{\nearrow}}
% \fr{..} : mot français (inventé par les modèles)
\providecommand{\fr}[1]{\textit{#1}}
\newtcolorbox{exemplebox}[1][]{popbase,colframe=poporange,before upper={\popboxhead{Exemple}{poporange}{#1}}}
\newtcolorbox{definition}[1][]{popbase,colframe=popblue,before upper={\popboxhead{Définition}{popblue}{#1}}}
\newtcolorbox{propriete}[1][]{popbase,colframe=poppurple,before upper={\popboxhead{Propriété}{poppurple}{#1}}}
\newtcolorbox{methode}[1][]{popbase,colframe=popgold,before upper={\popboxhead{Méthode}{popgold}{#1}}}
\newtcolorbox{attention}[1][]{popbase,colframe=poppink,before upper={\popboxhead{Attention}{poppink}{#1}}}
\newtcolorbox{experience}[1][]{popbase,colframe=popblue,colback=popblueL,before upper={\popboxhead{Expérience}{popblue}{#1}}}
% « Le savais-tu ? » : carte violette pleine (contexte, culture, Cameroun)
\newtcolorbox{savaistu}[1][]{popbase,colback=poppurple,colframe=popink,
  fontupper=\popsemi\fontsize{10.4}{14.2}\selectfont\color{white},
  before upper={\poppill{Le savais-tu\,?}{white}{poppurple}\ifx\relax#1\relax\else\hspace{6pt}{\popxbold\color{white}#1}\fi\par\vspace{7pt}}}
% Exercice résolu : énoncé en haut, \tcblower puis la solution sur fond crème
\newcounter{popexo}\newcounter{popexr}
\newtcolorbox{exoresolu}[1][]{popbase,colframe=poporange,colbacklower=popcream,
  segmentation style={popink,line width=1.2pt,dashed},
  before upper={\stepcounter{popexr}\popboxhead{Exercice résolu \thepopexr}{poporange}{#1}},
  before lower={\poppill{Solution}{popink}{popcream}\par\vspace{6pt}}}
% Exercice (non résolu en place) — la correction est dans la partie Corrigés
\newtcolorbox{exercice}[1][]{popbase,colframe=popink,
  before upper={\stepcounter{popexo}\popboxhead{Exercice \thepopexo}{popink}{#1}}}
% Corrigé
\newtcolorbox{corrige}[1][]{popbase,colframe=popgreen,colback=popgreenL,
  before upper={\popboxhead{Corrigé}{popgreen}{#1}}}
% Objectifs / prérequis / bilan
\newtcolorbox{objectifs}{popbase,colframe=popink,colback=poporangeL,
  before upper={\popboxhead{Objectifs de la leçon}{popink}{}}}
\newtcolorbox{prerequis}{popbase,colframe=popink,colback=popblueL,
  before upper={\popboxhead{Prérequis}{popblue}{}}}
\newtcolorbox{bilan}{popbase,colframe=popink,colback=white,boxrule=2.8pt,
  before upper={\popboxhead{Fiche bilan}{poporange}{L'essentiel à connaître pour l'examen}}}
% Figure encadrée avec légende
\newtcolorbox{popfigure}[1][]{enhanced,breakable=false,colback=white,colframe=popline,boxrule=1.4pt,arc=8pt,
  left=10pt,right=10pt,top=10pt,bottom=8pt,before skip=10pt,after skip=12pt,halign=center,
  lower separated=false,halign lower=center,
  fontlower=\popsemi\fontsize{9}{11.5}\selectfont\color{popmuted},
  #1}
\newcommand{\legende}[1]{\par\vspace{6pt}{\popsemi\fontsize{9}{11.5}\selectfont\color{popmuted}#1}}

% ============================================================
% Signature OnBuch+
% ============================================================
\newcommand{\poplogoinline}[1]{{\poplogo\fontsize{#1}{#1}\selectfont\color{popink}OnBuch\color{poporange}+}}
\newcommand{\signatureonbuch}{%
  \begin{tcolorbox}[enhanced,colback=popink,colframe=popink,boxrule=2.2pt,arc=10pt,
      drop shadow=poporange,left=14pt,right=14pt,top=10pt,bottom=10pt,before skip=0pt,after skip=0pt]
    \tikz[baseline=-0.7ex]\node[circle,fill=popgreen,draw=popcream,line width=1.3pt,minimum size=22pt,inner sep=0]{\popcheck{white}};%
    \hspace{9pt}\begin{minipage}[c]{0.62\linewidth}
      {\popxbold\fontsize{12}{13}\selectfont\color{popcream}Signé\ \textbullet\ {\poplogo OnBuch\color{poporange}+}}\par\vspace{2pt}
      {\raggedright\popsemi\fontsize{8}{10.5}\selectfont\color{popline}\MakeUppercase{Équipe pédagogique OnBuch+}\\\MakeUppercase{Conforme au programme officiel MINESEC}\par}
    \end{minipage}\hfill
    {\poplogo\fontsize{22}{22}\selectfont\color{popcream}OnBuch\color{poporange}+}
  \end{tcolorbox}%
}

% ============================================================
% Page de garde
% #1 titre · #2 matière · #3 niveau · #4 module · #5 n° leçon · #6 durée ·
% #7 description / objectifs (texte ou itemize)
% ============================================================
\newcommand{\pagegarde}[7]{%
  \thispagestyle{empty}
  \newgeometry{a4paper,margin=1.7cm,top=1.6cm,bottom=1.4cm}%
  \begin{tikzpicture}[remember picture,overlay]
    \fill[poporange!18] ([xshift=-3.2cm,yshift=-3.6cm]current page.north east) circle (4.6cm);
    \fill[poppurple!14] ([xshift=2.4cm,yshift=7.6cm]current page.south west) circle (3.8cm);
  \end{tikzpicture}%
  {\poplogo\fontsize{30}{30}\selectfont\color{popink}OnBuch\color{poporange}+}\hfill
  \raisebox{0.6ex}{\poppill{Cours signé \textbullet\ Premium}{popink}{popcream}}\par
  \vspace{1pt}\popsquiggle{poppurple}\par
  \vspace{8pt}
  \begin{tcolorbox}[enhanced,colback=poporange,colframe=popink,boxrule=2.8pt,arc=13pt,
      drop shadow=popink,left=18pt,right=18pt,top=12pt,bottom=13pt,after skip=10pt]
    \poppill{#2 \textbullet\ #3}{popink}{popcream}\hspace{5pt}\poppill{#4}{white}{popink}\par\vspace{8pt}
    {\popdisplay\fontsize{14}{14}\selectfont\color{popink}LEÇON #5}\par\vspace{4pt}
    {\raggedright\hyphenpenalty=10000\exhyphenpenalty=10000\popdisplay\fontsize{25}{27}\selectfont\color{popcream}\MakeUppercase{#1}\par}
    \vspace{8pt}
    {\popxbold\fontsize{9.5}{9.5}\selectfont\color{popink}\MakeUppercase{Durée conseillée : #6}}
  \end{tcolorbox}
  \begin{tcolorbox}[enhanced,colback=white,colframe=popink,boxrule=2.2pt,arc=10pt,
      drop shadow=popink,left=16pt,right=16pt,top=10pt,bottom=10pt,after skip=8pt]
    \poppill{Dans cette leçon}{poppurple}{white}\par\vspace{5pt}
    \popsemi\fontsize{9.3}{12.6}\selectfont\color{popink}#7
  \end{tcolorbox}
  \noindent\begin{minipage}[t]{0.487\linewidth}
    \begin{tcolorbox}[enhanced,colback=popgreen,colframe=popink,boxrule=2.2pt,arc=10pt,
        drop shadow=popink,left=11pt,right=11pt,top=10pt,bottom=10pt,height=3.1cm,valign=center]
      \tcbox[colback=white,colframe=popink,boxrule=1.3pt,arc=5pt,boxsep=0pt,nobeforeafter,
        left=3pt,right=3pt,top=3pt,bottom=3pt,valign=center]{\qrcode[height=1.9cm]{https://play.google.com/store/apps/details?id=com.luvvix.onbuch&hl=fr}}%
      \hspace{8pt}\begin{minipage}[c]{0.5\linewidth}
        {\popdisplay\fontsize{11}{12.5}\selectfont\color{white}\MakeUppercase{Télécharge OnBuch}}\par\vspace{4pt}
        {\raggedright\popsemi\fontsize{8.4}{11}\selectfont\color{white}Gratuit \textbullet\ Google Play\\Tuteur IA, annales, concours, résultats\par}
      \end{minipage}
    \end{tcolorbox}
  \end{minipage}\hfill
  \begin{minipage}[t]{0.487\linewidth}
    \begin{tcolorbox}[enhanced,colback=poppurple,colframe=popink,boxrule=2.2pt,arc=10pt,
        drop shadow=popink,left=11pt,right=11pt,top=10pt,bottom=10pt,height=3.1cm,valign=center]
      \tikz[baseline=-0.7ex]\node[circle,fill=white,draw=popink,line width=1.3pt,minimum size=1.6cm,inner sep=0]{\popdisplay\fontsize{24}{24}\selectfont\color{poppurple}+};%
      \hspace{8pt}\begin{minipage}[c]{0.6\linewidth}
        {\popdisplay\fontsize{11}{12.5}\selectfont\color{white}\MakeUppercase{Passe à OnBuch+}}\par\vspace{4pt}
        {\raggedright\popsemi\fontsize{8.4}{11}\selectfont\color{white}Tous les cours signés, Tuteur IA illimité, fiches PDF — onglet OnBuch+ de l'app\par}
      \end{minipage}
    \end{tcolorbox}
  \end{minipage}
  \vspace{8pt}
  \popcontenu
  \vfill
  \signatureonbuch
  \restoregeometry
  \newpage
}

% Rangée « Ce cours contient » (page de garde)
\newcommand{\poptile}[3]{% #1 couleur #2 gros texte #3 légende
  \begin{tcolorbox}[enhanced,colback=#1,colframe=popink,boxrule=2pt,arc=9pt,shadow={2.5pt}{-2.5pt}{0pt}{popink},
      left=6pt,right=6pt,top=9pt,bottom=9pt,halign=center,valign=center,height=1.75cm,width=\linewidth,nobeforeafter]
    {\popdisplay\fontsize{12.5}{13}\selectfont\color{white}\MakeUppercase{#2}}\par\vspace{3pt}
    {\centering\popsemi\fontsize{7.8}{9.5}\selectfont\color{white}#3\par}
  \end{tcolorbox}}
\newcommand{\popcontenu}{%
  \par\noindent{\popxboldls\fontsize{7.5}{7.5}\selectfont\color{popmuted}CE COURS SIGNÉ CONTIENT}\par\vspace{7pt}
  \noindent\begin{minipage}[t]{0.235\linewidth}\poptile{popblue}{Cours}{complet, illustré, pas à pas}\end{minipage}\hfill
  \begin{minipage}[t]{0.235\linewidth}\poptile{poporange}{Méthodes}{et exercices résolus}\end{minipage}\hfill
  \begin{minipage}[t]{0.235\linewidth}\poptile{poppink}{Exercices}{type examen + corrigés}\end{minipage}\hfill
  \begin{minipage}[t]{0.235\linewidth}\poptile{popgreen}{Fiche bilan}{l'essentiel à retenir}\end{minipage}\par}

% Sommaire
\newtcolorbox{sommaire}{enhanced,colback=white,colframe=popink,boxrule=2.2pt,arc=10pt,
  drop shadow=popink,left=18pt,right=18pt,top=13pt,bottom=13pt,before skip=6pt,after skip=20pt,
  before upper={\poppill{Sommaire}{poppurple}{white}\par\vspace{9pt}\popsemi\fontsize{10.6}{14.5}\selectfont\color{popink}}}

% Signature de fin de cours — poussée tout en bas de la dernière page
\newcommand{\findecours}{%
  \par\vfill\nopagebreak
  \begin{center}\popsquiggle{poporange}\end{center}\vspace{4pt}
  \signatureonbuch
}
\AtBeginDocument{\def\llbracket{\mathopen{[\mkern-3.5mu[}}\def\rrbracket{\mathclose{]\mkern-3.5mu]}}} % intervalles d'entiers (glyphe ⟦ absent de la police)
% Glyphes absents de Fira Math : redéfinis à partir de glyphes disponibles
\AtBeginDocument{%
  \def\setminus{\mathbin{\backslash}}\let\smallsetminus\setminus
  \def\vdots{\mathord{\vbox{\baselineskip3.2pt\lineskiplimit0pt\kern4pt\hbox{.}\hbox{.}\hbox{.}}}}%
  \def\ddots{\mathinner{\mkern1mu\raise6.5pt\hbox{.}\mkern2mu\raise3.6pt\hbox{.}\mkern2mu\raise0.7pt\hbox{.}\mkern1mu}}%
  \def\triangle{\mathord{\scalebox{0.9}{$\symup{\Delta}$}}}%
  \def\nabla{\mathord{\raisebox{\depth}{\scalebox{1}[-1]{$\symup{\Delta}$}}}}%
  \def\checkmark{\popcheck{popdark}}%
  \def\star{\mathbin{\ast}}\def\bigstar{\mathord{\ast}}%
  \def\diamond{\mathbin{\scalebox{0.6}{$\Diamond$}}}%
  \def\bigcup{\mathop{\mathchoice{\raisebox{-0.3em}{\scalebox{1.7}{$\cup$}}}{\scalebox{1.25}{$\cup$}}{\cup}{\cup}}\displaylimits}%
  \def\bigcap{\mathop{\mathchoice{\raisebox{-0.3em}{\scalebox{1.7}{$\cap$}}}{\scalebox{1.25}{$\cap$}}{\cap}{\cap}}\displaylimits}%
  \def\lesseqqgtr{\mathrel{\lessgtr}}\def\bigoplus{\mathop{\oplus}\displaylimits}%
}
\providecommand{\ssi}{si et seulement si} % abréviation fréquente
\AtBeginDocument{\providecommand{\wideparen}{\overparen}} % arc de cercle

\begin{document}

\pagegarde{Lesson 4 — Listening: Texts about taking part in communal work}{Anglais}{1ère A}{Chapter 1 : Family and Social Life}{ENA04}{2 h}{Cette leçon d'écoute entraîne les élèves à comprendre un texte lu ou un enregistrement sur la participation aux travaux communautaires. Elle développe les stratégies d'écoute (repérage du thème, des mots-clés, des chiffres et des noms propres), enrichit le lexique du travail collectif et prépare aux épreuves de compréhension orale du Bac.
\vspace{4pt}
\begin{multicols}{2}\small
\begin{itemize}
\item À la fin de cette leçon, je sais identifier le thème général d'un enregistrement dès la première écoute
\item À la fin de cette leçon, je sais repérer les mots-clés, les chiffres et les noms propres dans un texte écouté
\item À la fin de cette leçon, je sais prendre des notes efficaces pendant l'écoute
\item À la fin de cette leçon, je sais répondre à des questions de compréhension de types vrai/faux, QCM et réponses courtes
\item À la fin de cette leçon, je connais et j'emploie le lexique du travail communautaire (tools, tasks, organisation)
\item À la fin de cette leçon, je sais distinguer les informations essentielles des détails secondaires
\end{itemize}\end{multicols}}

\begin{sommaire}
\begin{multicols}{2}
\begin{enumerate}
\item Mise en route et découverte du thème
\item Script d'écoute : texte support
\item Compréhension : questions variées
\item Lexique thématique du travail communautaire
\item Méthode d'écoute et exemple traité
\item Production guidée : réemploi du lexique
\item Activité d'intégration
\item Méthodes et exercices résolus
\item Exercices
\item Corrigés des exercices
\item Fiche bilan
\end{enumerate}
\end{multicols}
\end{sommaire}

\begin{objectifs}
\begin{itemize}
\item À la fin de cette leçon, je sais identifier le thème général d'un enregistrement dès la première écoute
\item À la fin de cette leçon, je sais repérer les mots-clés, les chiffres et les noms propres dans un texte écouté
\item À la fin de cette leçon, je sais prendre des notes efficaces pendant l'écoute
\item À la fin de cette leçon, je sais répondre à des questions de compréhension de types vrai/faux, QCM et réponses courtes
\item À la fin de cette leçon, je connais et j'emploie le lexique du travail communautaire (tools, tasks, organisation)
\item À la fin de cette leçon, je sais distinguer les informations essentielles des détails secondaires
\item À la fin de cette leçon, je sais réemployer le lexique appris dans de courtes phrases personnelles
\item À la fin de cette leçon, je sais produire un court message oral ou écrit invitant à un travail communautaire
\end{itemize}
\end{objectifs}

\begin{prerequis}
\begin{itemize}
\item Le vocabulaire de base de la famille et de la vie sociale (Seconde, Unit 1)
\item Le présent simple et le présent continu pour décrire des habitudes et des actions
\item Les nombres, les dates et les heures en anglais
\item Les question words (who, what, where, when, why, how many)
\item Les prépositions de lieu et de temps (in, on, at, from... to...)
\end{itemize}
\end{prerequis}

\newpage

\coursec{Mise en route et découverte du thème}

C'est samedi matin à Bamenda. À sept heures, les haut-parleurs du quartier crépitent : \eng{“Everybody out! Today is our monthly community work!”} En quelques minutes, les rues se remplissent de voisins équipés de brouettes, de machettes et de balais. Les jeunes débouchent les caniveaux, les femmes balaient le marché, les hommes réparent le petit pont de bois sur la rivière. Personne n'est payé, mais tout le monde participe, car ce travail profite à tous. En Grande-Bretagne, on parlerait d'un \eng{volunteering day}, au Nigeria ou au Ghana de \eng{communal labour}. Au Cameroun anglophone, c'est tout simplement le \eng{community work}.

Voici maintenant la problématique de cette leçon : comment comprendre un reportage radio ou un discours en anglais qui décrit ce genre d'activité ? Quels mots-clés faut-il repérer ? Quelles informations faut-il noter ? C'est exactement ce que cette leçon d'écoute va vous apprendre. Mais avant d'écouter quoi que ce soit, un bon auditeur se prépare : il active ce qu'il sait déjà, il prédit le contenu, il se fixe des objectifs. C'est le travail de cette première section.

\courssub{Brainstorming : qu'est-ce que le \eng{communal work} ?}

Commençons par une question simple : que signifie \eng{communal work} ? Observez d'abord ces phrases, tirées de situations réelles de la vie camerounaise :

\begin{exemplebox}[Observer avant de définir]
\begin{itemize}
\item \eng{Every last Saturday of the month, the villagers clean the market together.} « Chaque dernier samedi du mois, les villageois nettoient le marché ensemble. »
\item \eng{The youth of the quarter built a wooden bridge over the stream.} « Les jeunes du quartier ont construit un pont en bois sur le ruisseau. »
\item \eng{During the dry season, people work on the chief's farm.} « Pendant la saison sèche, les gens travaillent dans le champ du chef. »
\item \eng{Nobody is paid for communal work, but everybody benefits from it.} « Personne n'est payé pour le travail communautaire, mais tout le monde en profite. »
\end{itemize}
\end{exemplebox}

Que remarquez-vous ? Dans chaque phrase, il y a un \cle{groupe de personnes} (\eng{the villagers}, \eng{the youth}, \eng{people}), une \cle{activité concrète} (\eng{clean}, \eng{built}, \eng{work}) et l'idée d'un \cle{bénéfice collectif} (\eng{everybody benefits}). Nous pouvons maintenant énoncer la définition.

\begin{definition}[Communal work]
\eng{Communal work} (ou \eng{community work}) désigne \eng{work done together by the members of a community for the benefit of all, usually without pay}, c'est-à-dire « un travail accompli ensemble par les membres d'une communauté, pour le bien de tous, généralement sans rémunération ».
\begin{itemize}
\item \eng{communal} (adjectif) : « communautaire, collectif » — de \eng{community}, « communauté ».
\item \eng{work} (nom indénombrable) : « travail ». On dit \eng{some work}, jamais \eng{a work} ni \eng{works} au sens de « travail ».
\end{itemize}
\end{definition}

\begin{attention}[Prononciation : COM-mu-nal]
L'adjectif \eng{communal} se prononce avec l'accent sur la \cle{première syllabe} : « KO-miou-nal ». Ne dites pas « commuNAL » sur le modèle français \emph{communal}. De même, \eng{community} s'accentue sur la deuxième syllabe : « ko-MIOU-ni-ti ». Une mauvaise accentuation peut empêcher un anglophone de vous comprendre.
\end{attention}

\courssub{Équivalences culturelles dans le monde anglophone}

Une même réalité sociale porte des noms différents selon les pays anglophones. Pour bien comprendre un enregistrement, il faut connaître ces variantes, car un reportage de la BBC n'emploiera pas le même mot qu'une radio de Bamenda ou de Lagos.

\begin{center}
\begin{tabularx}{\linewidth}{|X|X|X|}
\hline
\rowcolor{popblueL} Pays & Expression & Équivalent français et remarque \\
\hline
Cameroun anglophone & \eng{community work} & « travail communautaire » ; organisé à l'échelle du quartier ou du village \\
\hline
Nigeria, Ghana & \eng{communal labour} & « corvée collective » ; attention à l'orthographe britannique \eng{labour} (américain : \eng{labor}) \\
\hline
Grande-Bretagne, États-Unis & \eng{volunteering}, \eng{volunteer work} & « bénévolat » ; insiste sur le caractère volontaire \\
\hline
Grande-Bretagne & \eng{community service} & « service à la communauté » ; peut aussi désigner des travaux d'intérêt général \\
\hline
\end{tabularx}
\end{center}

\begin{attention}[Faux ami : \eng{volunteer} et « volontaire »]
\eng{To volunteer} signifie « se porter volontaire, faire du bénévolat » : \eng{She volunteers at the health centre.} « Elle est bénévole au centre de santé. » Mais attention : \eng{volunteer} signifie « bénévole », jamais « volontaire » au sens militaire (soldat volontaire) ou sportif. \eng{Voluntary work} se traduit par « travail bénévole » ou « bénévolat », non par « travail volontaire » (qui suggère « non forcé »). Retenez : \eng{volunteer} = « bénévole ».
\end{attention}

\begin{savaistu}[Le \eng{community work} au Cameroun anglophone]
Dans les régions du Nord-Ouest et du Sud-Ouest, le \eng{community work} est souvent organisé par le \eng{quarter head} (« chef de quartier ») ou par le \eng{traditional ruler} (« autorité traditionnelle », le \eng{fon}). La date est annoncée par les haut-parleurs ou par le \eng{town crier}. Et attention : ne pas y participer sans excuse valable peut entraîner une amende — \eng{a fine} ! Cette tradition renforce la solidarité : on y répare les routes, on nettoie les sources d'eau, on construit des salles de classe. Des pratiques semblables existent ailleurs : le \eng{harambee} au Kenya (« tirons ensemble ») ou le \eng{gotong royong} en Indonésie.
\end{savaistu}

\courssub{Carte mentale : organiser le vocabulaire du thème}

Pour préparer l'écoute, organisons nos idées autour de quatre questions : \eng{who?} (qui ?), \eng{what?} (quoi ?), \eng{where?} (où ?), \eng{when?} (quand ?). Cette carte mentale sera votre outil de prédiction : avant même d'entendre l'enregistrement, vous savez déjà quel type d'informations attendre.

\begin{popfigure}
\begin{tikzpicture}[mindmap, grow cyclic, every node/.style={concept}, root concept/.append style={concept color=popblue, font=\bfseries, minimum size=2.6cm, align=center}, level 1 concept/.append style={sibling angle=90, level distance=4cm, minimum size=2.2cm, font=\bfseries, align=center}, level 2 concept/.append style={minimum size=1.7cm, font=\small, align=center}]
\node[root concept]{COMMUNAL WORK}
  child[concept color=poporange] { node {WHO?}
    child { node {villagers} }
    child { node {youth} }
    child { node {leaders} }
  }
  child[concept color=popgreen] { node {WHAT?}
    child { node {cleaning} }
    child { node {building} }
    child { node {farming} }
  }
  child[concept color=poppurple] { node {WHERE?}
    child { node {market} }
    child { node {roads} }
    child { node {school} }
  }
  child[concept color=popgold] { node {WHEN?}
    child { node {Saturdays} }
    child { node {dry season} }
  };
\end{tikzpicture}
\legende{Carte mentale du thème \eng{communal work} : qui participe, quelles activités, quels lieux, quels moments.}
\end{popfigure}

\courssub{Activation du vocabulaire connu}

Vous connaissez déjà beaucoup de mots utiles depuis le collège. Les réactiver avant l'écoute facilite énormément la compréhension : le cerveau reconnaît plus vite un mot qu'il vient de « revoir ». Complétez mentalement ce tableau de réactivation.

\begin{center}
\begin{tabularx}{\linewidth}{|X|X|X|}
\hline
\rowcolor{popblueL} English & Français & Exemple en contexte \\
\hline
\eng{to clean} & nettoyer & \eng{They clean the market every month.} \\
\hline
\eng{to build} (built, built) & construire & \eng{The youth built a bridge.} \\
\hline
\eng{to help} & aider & \eng{Everyone helps the old people.} \\
\hline
\eng{the village} & le village & \eng{Our village is very united.} \\
\hline
\eng{people} (pluriel) & les gens & \eng{Many people take part in it.} \\
\hline
\eng{together} & ensemble & \eng{We work together.} \\
\hline
\eng{free of charge} & gratuitement & \eng{The work is done free of charge.} \\
\hline
\end{tabularx}
\end{center}

\begin{propriete}[Points de grammaire essentiels : \eng{people} et \eng{to build}]
\begin{itemize}
\item \textbf{\eng{people}} est déjà un \cle{pluriel} : on dit \eng{people are}, jamais \eng{peoples are} au sens de « les gens ». \eng{Peoples} existe mais signifie « les peuples » (groupes ethniques, nations).
\item \textbf{\eng{to build}} est un verbe irrégulier : \eng{build — built — built}.
\end{itemize}
\end{propriete}

\begin{center}
\begin{tabular}{|l|l|l|}
\hline
\rowcolor{popblueL} Forme & Anglais & Français \\
\hline
Infinitif & \eng{to build} & construire \\
\hline
Prétérit & \eng{built} & construisit / a construit \\
\hline
Participe passé & \eng{built} & construit \\
\hline
\end{tabular}
\end{center}

\courssub{Prédiction à partir du titre}

\begin{methode}[Prédire avant d'écouter : la stratégie du titre]
La prédiction est la première stratégie d'un bon auditeur. Voici la démarche en quatre étapes :
\begin{enumerate}
\item \textbf{Lire le titre} et souligner les mots porteurs de sens : ici, \eng{taking part in communal work}.
\item \textbf{Traduire mentalement} : « participer au travail communautaire ». Le verbe \eng{to take part in} signifie « participer à ».
\item \textbf{Se poser les questions-clés} : \eng{who? what? where? when? why?} (qui, quoi, où, quand, pourquoi).
\item \textbf{Formuler des hypothèses} : « L'enregistrement parlera probablement de villageois qui nettoient ou construisent quelque chose, un samedi, dans une localité précise, avec peut-être des chiffres et des noms propres. »
\end{enumerate}
\end{methode}

Appliquons cette méthode. Le titre de notre leçon est \eng{Taking part in communal work}. Le mot \eng{taking part} nous annonce qu'on parlera de \cle{participation} : on s'attend donc à entendre qui participe (des habitants, des jeunes, des chefs), et peut-être qui refuse de participer. Le mot \eng{communal work} annonce des activités concrètes : nettoyage, construction, travaux agricoles. Un reportage radio contiendra aussi très probablement des \cle{noms propres} (un village, une école), des \cle{chiffres} (le nombre de participants, une heure, une date) et peut-être une \cle{interview} d'un participant ou d'un responsable.

\begin{exoresolu}[Prédire le contenu d'un enregistrement]
Énoncé : le titre d'un reportage radio est \eng{“Buea villagers clean their water source”}. Avant l'écoute, formulez trois prédictions sur ce que vous allez entendre, en utilisant \eng{who}, \eng{what}, \eng{where}.
\tcblower
Solution détaillée : on applique la stratégie du titre, étape par étape.
\begin{enumerate}
\item Mots porteurs de sens : \eng{villagers} (qui), \eng{clean} (quoi), \eng{water source} (où/quoi), \eng{Buea} (où).
\item Prédiction \eng{who} : \eng{The speakers will probably be villagers of Buea, maybe a quarter head and a young participant.} « Les intervenants seront probablement des villageois de Buea, peut-être un chef de quartier et un jeune participant. »
\item Prédiction \eng{what} : \eng{They will talk about cleaning the water source: removing dirt, cutting grass around it.} « Ils parleront du nettoyage de la source : enlever les saletés, couper l'herbe autour. »
\item Prédiction \eng{where/when} : \eng{The report may mention the exact place and the day of the activity, perhaps a Saturday.} « Le reportage mentionnera peut-être le lieu exact et le jour de l'activité, probablement un samedi. »
\end{enumerate}
Ces prédictions orientent l'écoute : au lieu de tout comprendre mot à mot, on vérifie ses hypothèses et on repère les informations attendues.
\end{exoresolu}

\courssub{Un premier texte pour entrer dans le thème}

Avant l'écoute proprement dite (section suivante), lisez ce court texte d'entraînement. Il utilise exactement le type de langue que vous entendrez dans l'enregistrement.

\begin{textebox}[A radio announcement in Bamenda]
\eng{Good morning, dear listeners! This is your community radio speaking. As you all know, next Saturday is our monthly community work day. All the inhabitants of Nkwen quarter are invited to take part. We shall meet at the market square at seven o'clock sharp. Please come with your own tools: brooms, machetes, hoes and wheelbarrows.}

\eng{This month, our main task is to clean the market and to unblock the gutters before the rainy season. The youth will also repair the wooden bridge near the primary school. Our quarter head, Mr. Ngong, reminds everybody that community work is not optional: anyone who stays at home without a good reason will pay a fine of two thousand francs.}

\eng{Last month, more than three hundred people took part, and together we cleaned the whole market in only four hours. The women swept the stalls, the men emptied the dustbins, and the children picked up the plastic bags. It was hard work, but it was also a pleasure: we sang, we laughed, and at the end, the women's association offered food and palm wine to everybody.}

\eng{So, dear listeners, do not stay in bed next Saturday! Community work keeps our quarter clean, safe and united. Remember: a clean quarter is a healthy quarter. See you on Saturday at seven o'clock, at the market square!}
\end{textebox}

\textbf{Glossaire}

\begin{center}
\begin{tabularx}{\linewidth}{|X|X|X|}
\hline
\rowcolor{popblueL} Mot anglais & Nature & Traduction \\
\hline
\eng{tools} & nom pluriel & outils \\
\hline
\eng{broom} & nom & balai \\
\hline
\eng{hoe} & nom & houe \\
\hline
\eng{wheelbarrow} & nom & brouette \\
\hline
\eng{task} & nom & tâche \\
\hline
\eng{to unblock the gutters} & expression & déboucher les caniveaux \\
\hline
\eng{optional} & adjectif & facultatif \\
\hline
\eng{fine} & nom & amende \\
\hline
\eng{stall} & nom & étal (de marché) \\
\hline
\eng{to sweep} (swept, swept) & verbe irrégulier & balayer \\
\hline
\eng{sharp} & adverbe & (à sept heures) précises \\
\hline
\end{tabularx}
\end{center}

\textbf{Questions de pré-compréhension} (répondez mentalement, les réponses serviront dans la section suivante) :

\begin{enumerate}
\item \eng{When will the community work take place?} « Quand le travail communautaire aura-t-il lieu ? »
\item \eng{What are the two main tasks this month?} « Quelles sont les deux tâches principales ce mois-ci ? »
\item \eng{What happens to people who do not participate?} « Qu'arrive-t-il aux personnes qui ne participent pas ? »
\end{enumerate}

\courssub{Les objectifs d'écoute : du sens global aux détails}

Terminons cette mise en route en fixant clairement nos objectifs. Une bonne écoute se fait toujours en \cle{deux passages}.

\begin{aretenir}[Écouter en deux temps]
\begin{itemize}
\item \textbf{Première écoute — sens global} (\eng{listening for gist}) : ne notez rien, écoutez simplement pour répondre à trois questions : \emph{de quoi parle-t-on ? qui parle ? où et quand cela se passe-t-il ?} Ne paniquez pas si vous ne comprenez pas chaque mot : c'est normal et inutile.
\item \textbf{Deuxième écoute — les détails} (\eng{listening for detail}) : cette fois, notez les informations précises : les \cle{chiffres} (dates, heures, nombres de participants, montants), les \cle{noms propres} (personnes, lieux) et les \cle{mots-clés} du thème (\eng{clean, build, repair, fine, tools}).
\end{itemize}
Règle d'or : on ne cherche jamais à traduire mot à mot en écoutant. On écoute pour \emph{repérer}, pas pour \emph{traduire}.
\end{aretenir}

\begin{experience}[Brainstorming collectif en classe]
Par groupes de quatre, en cinq minutes :
\begin{enumerate}
\item Chaque groupe dresse la liste de tous les exemples de \eng{communal work} qu'il connaît dans sa ville ou son village (nettoyage du quartier, construction d'un pont, travaux au champ du chef, tournée d'eau à l'école).
\item Chaque groupe choisit un exemple et le décrit en deux phrases anglaises simples, par exemple : \eng{In our quarter, people clean the road every Saturday. The youth carry the sand.} « Dans notre quartier, les gens nettoient la route chaque samedi. Les jeunes transportent le sable. »
\item Un porte-parole par groupe présente ses phrases à la classe ; le professeur note au tableau le nouveau vocabulaire qui apparaît.
\end{enumerate}
Cette activité réactive vos connaissances et enrichit la carte mentale avant l'écoute.
\end{experience}

Vous êtes maintenant prêts : le thème est clair, le vocabulaire de base est activé, vos prédictions sont faites et vos objectifs d'écoute sont fixés. Passons au script d'écoute lui-même.

\coursec{Script d'écoute : texte support}

Dans la section précédente, nous avons découvert le thème du \cle{communal work} (travail communautaire) et formulé des prédictions : qui participe ? quelles tâches sont accomplies ? où et quand ? Le moment est venu de vérifier ces hypothèses grâce à un document sonore authentique : un reportage de la radio locale de Bafut, petite ville du Nord-Ouest du Cameroun. L'enseignant lit le script \textbf{deux fois} ; vous n'écrivez rien pendant la première écoute, vous prenez des notes pendant la seconde.

\courssub{Première écoute : saisir le sens global}

\begin{methode}[Première écoute : les quatre questions]
Avant d'appuyer sur « play » (ou avant que l'enseignant ne lise), fixez-vous un objectif simple : répondre à quatre questions, sans chercher les détails.
\begin{enumerate}
\item \textbf{Who?} — Qui parle ? De qui parle-t-on ? (un présentateur, un chef de quartier, des villageois)
\item \textbf{What?} — De quoi s'agit-il ? (une journée de travail communautaire)
\item \textbf{Where?} — Où cela se passe-t-il ? (à Bafut, autour du marché, près de l'école)
\item \textbf{When?} — Quand ? (samedi dernier, de 7 h à midi)
\end{enumerate}
Pendant l'écoute, ne vous accrochez pas à un mot inconnu : laissez-le passer et captez les mots que vous reconnaissez (\eng{village}, \eng{market}, \eng{bridge}, \eng{work}). Le sens global naît de la combinaison de ces mots familiers.
\end{methode}

\begin{attention}[Ne paniquez pas sur les mots inconnus]
L'erreur classique de l'élève francophone est de « bloquer » sur un mot comme \eng{gutters} ou \eng{wheelbarrows} et de perdre le fil de tout le passage. En anglais comme en français, on comprend un message sans comprendre chaque mot. Entraînez-vous à tolérer le flou : la deuxième écoute et le glossaire viendront combler les lacunes.
\end{attention}

\courssub{Le script : « Community Work in Bafut »}

\begin{textebox}[Radio report — “Community Work in Bafut”]
Good morning, listeners! This is Radio Bafut. Last Saturday, more than two hundred villagers took part in communal work in Bafut. The quarter head, Mr Ngum, told us: “Every able-bodied man and woman must join.”

From 7 a.m. to noon, the youth cleared the gutters around the market, while the women swept the streets and the men repaired the wooden bridge near the primary school. The council provided wheelbarrows, spades and cutlasses.

“This work keeps our village clean and united,” said Mrs Abong, a trader.

Next month, the community plans to dig a new well behind the health centre. So, dear listeners, come and join us — together we are stronger!
\end{textebox}

\textbf{Traduction guidée (à cacher pendant l'écoute).} « Bonjour chers auditeurs ! Ici Radio Bafut. Samedi dernier, plus de deux cents villageois ont participé au travail communautaire à Bafut. Le chef de quartier, M. Ngum, nous a dit : “Tout homme et toute femme valide doit se joindre à nous.” De 7 heures à midi, les jeunes ont débouché les caniveaux autour du marché, pendant que les femmes balayaient les rues et que les hommes réparaient le pont en bois près de l'école primaire. La commune a fourni des brouettes, des pelles et des machettes. “Ce travail garde notre village propre et uni”, a déclaré Mme Abong, une commerçante. Le mois prochain, la communauté prévoit de creuser un nouveau puits derrière le centre de santé. Alors, chers auditeurs, venez vous joindre à nous — ensemble nous sommes plus forts ! »

\courssub{Deuxième écoute : la prise de notes ciblée}

\begin{methode}[Deuxième écoute : noter les détails]
La deuxième écoute vise les \cle{details}. Préparez une grille à trois colonnes et remplissez-la pendant l'écoute :
\begin{enumerate}
\item \textbf{Numbers} (chiffres) : combien de participants ? quelles heures ?
\item \textbf{Proper names} (noms propres) : personnes, lieux, institutions.
\item \textbf{Tasks} (tâches) : qui fait quoi ?
\end{enumerate}
Notez des mots, pas des phrases complètes : \eng{200 villagers / 7 a.m.–noon / Mr Ngum / youth → gutters}. Vous rédigerez des phrases entières après l'écoute.
\end{methode}

\begin{center}
\begin{tabularx}{\linewidth}{|X|X|X|}
\hline
\rowcolor{popblueL} \textbf{Numbers} & \textbf{Proper names} & \textbf{Tasks} \\
\hline
two hundred villagers & Bafut (village) & the youth cleared the gutters \\
\hline
from 7 a.m. to noon & Mr Ngum (quarter head) & the women swept the streets \\
\hline
next month (new well) & Mrs Abong (trader) & the men repaired the bridge \\
\hline
 & the council & the council provided tools \\
\hline
\end{tabularx}
\end{center}

\begin{propriete}[Structure d'un reportage radio]
Un reportage radio anglophone suit presque toujours la même architecture : le présentateur \textbf{introduit} le sujet, une personne interviewée apporte un \textbf{témoignage}, le reporter donne les \textbf{détails} (chiffres, tâches, lieux), puis une \textbf{conclusion} appelle à l'action ou annonce la suite. Repérer cette structure pendant l'écoute aide à anticiper le type d'information qui va suivre.
\end{propriete}

\begin{popfigure}
\begin{tikzpicture}[
  node distance=0.45cm,
  box/.style={draw=popblue, thick, rounded corners=4pt, fill=popblueL, align=center, text width=4.6cm, minimum height=1.1cm, font=\small},
  lbl/.style={font=\small\itshape, text=popmuted, align=center, text width=3.2cm},
  arr/.style={-{Stealth[length=3mm]}, thick, poporange}
]
\node[box, align=center] (intro){\textbf{Introduction}\\ presenter: “Good morning, listeners!”};
\node[box, below=of intro, align=center] (interview){\textbf{Interview}\\ quarter head: “Every able-bodied man and woman must join.”};
\node[box, below=of interview, align=center] (details){\textbf{Details}\\ tasks, numbers: 200 villagers, 7 a.m.–noon};
\node[box, below=of details, align=center] (concl){\textbf{Conclusion}\\ appeal: “Come and join us!”};
\node[lbl, right=0.4cm of intro] {who, what, where, when};
\node[lbl, right=0.4cm of interview] {a voice from the community};
\node[lbl, right=0.4cm of details] {facts and figures};
\node[lbl, right=0.4cm of concl] {call to action};
\draw[arr] (intro) -- (interview);
\draw[arr] (interview) -- (details);
\draw[arr] (details) -- (concl);
\end{tikzpicture}
\legende{Structure du reportage « Community Work in Bafut » : de l'introduction à l'appel final.}
\end{popfigure}

\courssub{Vérification collective des prédictions}

Après les deux écoutes, la classe confronte ses prédictions (section 1) au contenu réel du script. C'est une étape essentielle : elle montre que l'on peut anticiper le contenu d'un document sonore grâce au titre et au contexte.

\begin{center}
\begin{tabularx}{\linewidth}{|X|X|X|}
\hline
\rowcolor{popblueL} \textbf{Prédiction possible} & \textbf{Réalité du script} & \textbf{Verdict} \\
\hline
People will clean the village. & They cleared gutters, swept streets, repaired a bridge. & Confirmé (et précisé) \\
\hline
A chief will speak. & The quarter head, Mr Ngum, speaks. & Confirmé \\
\hline
The work will last one day. & From 7 a.m. to noon (half a day). & Partiellement confirmé \\
\hline
They will build a school. & They plan to dig a well next month. & Infirmé \\
\hline
\end{tabularx}
\end{center}

\begin{exemplebox}[Phrases de vérification à l'oral]
\eng{We predicted that people would clean the village, and we were right.} « Nous avions prédit que les gens nettoieraient le village, et nous avions raison. »

\eng{We thought they would build a school, but in fact they plan to dig a well.} « Nous pensions qu'ils construiraient une école, mais en fait ils prévoient de creuser un puits. »

\eng{The work did not last the whole day: it lasted from 7 a.m. to noon.} « Le travail n'a pas duré toute la journée : il a duré de 7 h à midi. »
\end{exemplebox}

\courssub{Lecture silencieuse du script et glossaire}

Après l'écoute, lisez le script en silence. Cette lecture confirme ce que l'oreille a saisi et révèle les mots que vous aviez mal identifiés (par exemple \eng{spades}, souvent confondu avec \eng{spoons} à l'écoute).

\begin{definition}[Glossaire du reportage]
\begin{itemize}
\item \cle{able-bodied} (adj.) [« é-ble bo-di »] : physically strong and healthy enough to work — valide, en bonne santé, apte au travail. \eng{Every able-bodied man and woman must join.} « Tout homme et toute femme valide doit participer. »
\item \cle{quarter head} (n.) [« kwor-tor hed »] : the leader of a neighbourhood — chef de quartier.
\item \cle{gutter} (n.) [« ga-teur »] : a channel at the side of a road that carries away rainwater — caniveau.
\item \cle{to sweep} (v., swept, swept) [« sou-ip », « suèpte »] : balayer.
\item \cle{wheelbarrow} (n.) [« wil-ba-rou »] : brouette.
\item \cle{spade} (n.) [« speïd »] : pelle.
\item \cle{cutlass} (n.) [« cat-lass »] : machette.
\item \cle{trader} (n.) [« trè-deur »] : commerçant(e), marchand(e).
\item \cle{to dig} (v., dug, dug) [« dig », « dag »] : creuser.
\item \cle{well} (n.) [« ouèl »] : puits.
\item \cle{council} (n.) [« koun-sill »] : la commune, le conseil municipal.
\item \cle{health centre} (n.) [« èlte senteur »] : centre de santé.
\end{itemize}
\end{definition}

\begin{exemplebox}[Exemples traduits]
\eng{The youth cleared the gutters.} « Les jeunes ont débouché les caniveaux. »

\eng{The council provided wheelbarrows, spades and cutlasses.} « La commune a fourni des brouettes, des pelles et des machettes. »

\eng{The men repaired the wooden bridge near the primary school.} « Les hommes ont réparé le pont en bois près de l'école primaire. »

\eng{This work keeps our village clean and united.} « Ce travail garde notre village propre et uni. »
\end{exemplebox}

\begin{attention}[« Two hundred » ne prend jamais de -s]
En anglais, \cle{hundred}, \cle{thousand} et \cle{million} restent \textbf{invariables} quand ils sont précédés d'un chiffre ou de \eng{a} :
\begin{itemize}
\item \eng{two hundred villagers} « deux cents villageois » — jamais \emph{two hundreds}.
\item \eng{three thousand people} « trois mille personnes ».
\item \eng{a million inhabitants} « un million d'habitants ».
\end{itemize}
Le -s n'apparaît que dans l'expression indéfinie \eng{hundreds of people} « des centaines de personnes », \eng{thousands of villagers} « des milliers de villageois ». Attention au contraste avec le français : « deux \textbf{cents} » prend un -s, « two hundred » non !
\end{attention}

\begin{exoresolu}[Repérage d'informations dans le script]
Énoncé — Répondez en phrases complètes, en anglais, à partir du script : (1) How many villagers took part in the communal work? (2) Who is Mr Ngum? (3) What did the women do? (4) What tools did the council provide? (5) What does the community plan to do next month?
\tcblower
Solution détaillée :
\begin{enumerate}
\item \eng{More than two hundred villagers took part in the communal work.} — On repère le chiffre dans la première phrase du reportage (« more than two hundred villagers »).
\item \eng{Mr Ngum is the quarter head.} — Le script le présente par une apposition : « The quarter head, Mr Ngum ».
\item \eng{The women swept the streets.} — Attention au parallélisme du passage : the youth \emph{cleared}, the women \emph{swept}, the men \emph{repaired}.
\item \eng{The council provided wheelbarrows, spades and cutlasses.} — Les trois outils sont énumérés dans la même phrase.
\item \eng{Next month, the community plans to dig a new well behind the health centre.} — L'information se trouve dans la conclusion, après le témoignage de Mrs Abong.
\end{enumerate}
Méthode : pour chaque question, repérez d'abord le \textbf{mot-clé} de la question (\eng{how many}, \eng{who}, \eng{what}), puis localisez ce mot (ou son équivalent) dans le script, et recopiez la phrase en l'adaptant.
\end{exoresolu}

\begin{savaistu}[Le travail communautaire au Cameroun anglophone]
Dans les régions du Nord-Ouest et du Sud-Ouest du Cameroun, le travail communautaire est une tradition vivante : les habitants d'un village ou d'un quartier se réunissent régulièrement pour nettoyer les rues, déboucher les caniveaux avant la saison des pluies, réparer les ponts ou construire des puits. En anglais, on parle de \eng{communal work} ou de \eng{community work}. Les radios locales comme « Radio Bafut » jouent un rôle clé : elles annoncent les dates, mobilisent les habitants et félicitent les participants. L'expression finale du reportage, \eng{together we are stronger} (« ensemble nous sommes plus forts »), résume l'esprit de solidarité qui anime ces journées.
\end{savaistu}

\courssub{Production guidée : réemploi du lexique}

\begin{experience}[Jeu de rôle : « Radio Bafut » — Expression orale]
En binôme, préparez une courte interview radio (1 à 2 minutes) à partir du reportage.
\begin{enumerate}
\item \textbf{Rôle A (Journaliste) :} Vous êtes l'animateur de Radio Bafut. Posez des questions pour obtenir des détails : \eng{How many people came? What did the women do? What tools did the council give? What is the next project?}
\item \textbf{Rôle B (Témoin / Chef de quartier) :} Vous avez participé. Répondez en réemployant le vocabulaire : \eng{able-bodied, gutters, swept, wheelbarrows, spades, cutlasses, well, health centre}.
\end{enumerate}
Changez de rôle. L'enseignant circule et note la prononciation de \eng{gutters} [« ga-teur »], \eng{wheelbarrows} [« wil-ba-rouz »], \eng{spades} [« speïdz »], \eng{cutlasses} [« cat-lass-iz »].
\end{experience}

\begin{exercice}[À vous de jouer — Complétion]
Écoutez une dernière fois le script (ou relisez-le), puis complétez les phrases suivantes avec le mot qui convient : \emph{gutters, wheelbarrows, quarter head, well, swept, able-bodied}.
\begin{enumerate}
\item Every \ldots\ldots\ldots man and woman must join the work.
\item The youth cleared the \ldots\ldots\ldots around the market.
\item The women \ldots\ldots\ldots the streets.
\item The council provided \ldots\ldots\ldots, spades and cutlasses.
\item The community plans to dig a new \ldots\ldots\ldots behind the health centre.
\item Mr Ngum is the \ldots\ldots\ldots of Bafut.
\end{enumerate}
\end{exercice}

\begin{corrige}[Exercice « À vous de jouer »]
1. able-bodied — 2. gutters — 3. swept — 4. wheelbarrows — 5. well — 6. quarter head. Astuce : chaque réponse se vérifie en retrouvant la phrase exacte dans le script ; relisez la phrase complète à voix haute pour mémoriser le mot dans son contexte.
\end{corrige}

\begin{exercice}[Production écrite courte — Mini-reportage]
Rédigez un paragraphe de 5 à 6 lignes en anglais pour le journal de l'école. Titre : \eng{“Community Work in Our Neighbourhood”}. Utilisez au moins 6 mots du glossaire (\eng{able-bodied, quarter head, gutter, sweep, wheelbarrow, spade, cutlass, trader, dig, well, council, health centre}). Respectez la structure : introduction → détails (chiffres, tâches) → citation → projet futur → appel.
\end{exercice}

\begin{corrige}[Exercice « Production écrite » — Proposition de correction]
\eng{“Community Work in Our Neighbourhood”}

\eng{Last Saturday, more than one hundred able-bodied villagers took part in communal work in our quarter. The quarter head, Mr Bello, encouraged everyone to join. From 7 a.m. to 11 a.m., the youth cleared the gutters near the market while the women swept the streets. The men repaired the bridge with spades and cutlasses. The council provided wheelbarrows. “This work unites us,” said Mrs Atanga, a trader. Next month, we plan to dig a well behind the health centre. Come and join us — together we are stronger!}

\textbf{Analyse :} Le texte reprend la structure du modèle (introduction, détails parallèles, citation, projet futur, appel). Le vocabulaire ciblé est intégré naturellement. Les temps (past simple pour le récit, present simple pour la citation et le projet futur avec \eng{plan to}) sont corrects.
\end{corrige}

Vous savez désormais extraire le sens global puis les détails d'un reportage radio, et réemployer le lexique du travail communautaire à l'oral comme à l'écrit. La section suivante vous entraînera à répondre à des questions de compréhension variées : vrai/faux, QCM et réponses courtes.

\coursec{Compréhension : questions variées}

Vous venez d'écouter (ou de lire) le script sur les travaux communautaires du village de Mbankomo. Pour vérifier que vous avez bien compris l'enregistrement, l'examinateur — comme au probatoire et au baccalauréat — vous propose plusieurs types de questions : des affirmations à juger \eng{true} ou \eng{false}, des questions à choix multiples (QCM), des questions à réponses courtes et une fiche d'informations à compléter. Chaque type de question obéit à une méthode précise. C'est ce que nous allons apprendre pas à pas dans cette section.

\begin{textebox}[Script d'écoute — “Community Work in Mbankomo”]
Last Saturday, the people of Mbankomo took part in communal work. The quarter head, Mr. Etoundi, had announced the work on the community radio two weeks earlier. More than two hundred villagers answered the call. The work started at seven o'clock in the morning and finished at noon.

The men cleared the road to the market with machetes and hoes, while the women swept the market square and collected the rubbish in baskets. The young boys carried water, and the girls helped to plant flowers around the health centre. The traders of the market offered food and palm wine to the workers at break time.

At the end of the day, the quarter head thanked everybody. “A clean village is a healthy village,” he said. He also announced that next month, the community will dig a new well near the primary school, because the old one dries up during the dry season. Everybody promised to come again.
\end{textebox}

\begin{definition}[Glossaire : mots-clés du texte]
\begin{itemize}
\item \eng{quarter head} (n.) [« kwɔ-tə hed »] : chef de quartier
\item \eng{to clear} (v.) [« klɪə »] : débroussailler, nettoyer
\item \eng{hoe} (n.) [« həʊ »] : houe
\item \eng{machete} (n.) [« mə-ʃe-ti »] : machette
\item \eng{rubbish} (n.) [« rʌ-bɪʃ »] : ordures (anglais britannique)
\item \eng{break time} (n.) [« breɪk taɪm »] : pause, moment de repos
\item \eng{well} (n.) [« wel »] : puits
\item \eng{to dry up} (v.) [« draɪ ʌp »] : s'assécher, tarir
\item \eng{palm wine} (n.) [« pɑːm waɪn »] : vin de palme
\end{itemize}
\end{definition}

\courssub{Le Vrai / Faux (\eng{True or False}) : observer avant de juger}

\begin{methode}[Répondre à un Vrai / Faux en quatre étapes]
\begin{enumerate}
\item \textbf{Avant l'écoute}, lisez chaque affirmation et soulignez les mots-clés : les \cle{dates}, les \cle{chiffres}, les \cle{noms propres}, les \cle{verbes d'action}.
\item \textbf{Pendant l'écoute}, notez en face de chaque mot-clé l'information entendue (un mot suffit : \eng{Saturday}, \eng{200}, \eng{well}).
\item \textbf{Comparez} l'affirmation avec vos notes : si un seul détail est différent, l'affirmation est \eng{false}.
\item \textbf{Justifiez toujours} votre réponse en citant la phrase exacte du texte qui prouve votre choix.
\end{enumerate}
\end{methode}

\begin{exemplebox}[Exemple traité pas à pas]
Affirmation : \eng{The work took place on Sunday.}

Étape 1 — mots-clés soulignés : \eng{took place} + \eng{Sunday} (jour).

Étape 2 — pendant l'écoute, on entend : \eng{Last Saturday, the people of Mbankomo took part in communal work.}

Étape 3 — comparaison : \eng{Sunday} $\neq$ \eng{Saturday} : l'affirmation est fausse.

Étape 4 — réponse rédigée : \eng{False. The text says: “Last Saturday, the people of Mbankomo took part in communal work.”} « Faux. Le texte dit : “Samedi dernier, les habitants de Mbankomo ont participé aux travaux communautaires.” »
\end{exemplebox}

\begin{propriete}[La justification est obligatoire]
Au baccalauréat et au probatoire, une réponse \eng{True} ou \eng{False} juste \textbf{mais non justifiée} peut perdre la moitié des points. La justification se fait en \textbf{citant la phrase du texte} (entre guillemets anglais) ou en la paraphrasant brièvement. Formule type : \eng{True / False. The text says: “...”} ou \eng{False, because the speaker says that...}
\end{propriete}

\courssub{Les questions à choix multiples (QCM)}

Dans un QCM, une seule réponse est correcte ; les autres sont des \cle{leurres} (\eng{distractors}) : ce sont souvent des mots qui apparaissent réellement dans le texte, mais qui ne répondent pas à la question. C'est exactement le piège tendu ici.

\begin{exoresolu}[QCM : Who organised the work?]
Énoncé : \eng{Who organised the work? a) the mayor \quad b) the quarter head \quad c) the traders}
\tcblower
Solution : la bonne réponse est \textbf{b) the quarter head}.

Justification : \eng{The quarter head, Mr. Etoundi, had announced the work on the community radio.} « Le chef de quartier, M. Etoundi, avait annoncé les travaux à la radio communautaire. »

Pourquoi les autres réponses sont des leurres : \eng{the traders} apparaît bien dans le texte, mais ils ont seulement \emph{offert} la nourriture (\eng{offered food}) ; \eng{the mayor} n'est jamais mentionné. Un mot entendu n'est pas forcément la bonne réponse : il faut vérifier \emph{qui fait quoi}.
\end{exoresolu}

\begin{attention}[Piège du QCM : le mot entendu n'est pas la réponse]
Les concepteurs d'épreuves placent dans les mauvaises options des mots du texte pour piéger les élèves qui écoutent « par mots » au lieu d'écouter le sens. Avant de cocher, demandez-vous : \emph{ce mot répond-il vraiment à la question posée ?} Ici, la question porte sur \eng{who organised} (qui a organisé), pas sur \eng{who offered food} (qui a offert à manger).
\end{attention}

\courssub{Les questions à réponses courtes (\eng{short-answer questions})}

Les questions ouvertes commencent par un mot interrogatif : \eng{How many...?} (combien de... ?), \eng{What...?} (que / quoi... ?), \eng{When...?} (quand... ?), \eng{Where...?} (où... ?), \eng{Who...?} (qui... ?), \eng{Why...?} (pourquoi... ?). La réponse doit être \textbf{directe et complète grammaticalement}, sans recopier la question.

\begin{exemplebox}[Exemples traduits]
\eng{How many villagers took part?} « Combien de villageois ont participé ? »

Réponse : \eng{More than two hundred (villagers took part).} « Plus de deux cents (villageois ont participé). »

\eng{What will the community dig next month?} « Que creusera la communauté le mois prochain ? »

Réponse : \eng{A new well near the primary school.} « Un nouveau puits, près de l'école primaire. »

\eng{When did the work start and finish?} « Quand les travaux ont-ils commencé et fini ? »

Réponse : \eng{From 7 a.m. to noon.} « De 7 heures du matin à midi. »
\end{exemplebox}

\begin{attention}[Ne recopiez pas toute la question]
Erreur fréquente des francophones : répéter la question dans la réponse, ce qui produit des phrases redondantes ou incorrectes.

Mauvais : \eng{The work lasted from 7 a.m. to noon it lasted.}

Correct : \eng{From 7 a.m. to noon.} ou, en phrase complète : \eng{The work lasted from 7 a.m. to noon.}

Règle d'or : répondez \textbf{directement} à la question, comme en conversation. La réponse courte est acceptée si elle contient l'information exacte et une grammaire correcte.
\end{attention}

\begin{center}
\begin{tabularx}{\linewidth}{|X|X|X|}
\hline
\rowcolor{popblueL} Mot interrogatif & Question du texte & Réponse attendue \\
\hline
\eng{When} (quand) & \eng{When did the work take place?} & \eng{Last Saturday.} \\
\hline
\eng{Where} (où) & \eng{Where did the women work?} & \eng{In the market square.} \\
\hline
\eng{Who} (qui) & \eng{Who offered food to the workers?} & \eng{The traders.} \\
\hline
\eng{How many} (combien) & \eng{How many villagers took part?} & \eng{More than two hundred.} \\
\hline
\eng{What} (quoi) & \eng{What will they dig next month?} & \eng{A new well.} \\
\hline
\eng{Why} (pourquoi) & \eng{Why will they dig a new well?} & \eng{Because the old one dries up.} \\
\hline
\end{tabularx}
\end{center}

Remarquez que \eng{how many} est suivi d'un nom \textbf{pluriel} (\eng{how many villagers}), alors que le français « combien de » accepte les deux. Et notez que \eng{why} appelle une réponse introduite par \eng{because} : \eng{Because the old well dries up during the dry season.} « Parce que l'ancien puits s'assèche pendant la saison sèche. »

\courssub{Lexique thématique : le travail communautaire}

Pour parler des travaux communautaires dans votre quartier ou votre village, vous devez maîtriser ce vocabulaire de base, classé par catégories.

\begin{definition}[Vocabulaire essentiel — Communal Work]
\begin{center}
\begin{tabularx}{\linewidth}{|l|X|l|}
\hline
\rowcolor{popblueL} \textbf{Personnes / Rôles} & \textbf{Actions / Verbes} & \textbf{Outils / Matériel} \\
\hline
\eng{quarter head} (chef de quartier) & \eng{to clear} (débroussailler) & \eng{machete} (machette) \\
\eng{villagers} (villageois) & \eng{to sweep} (balayer) & \eng{hoe} (houe) \\
\eng{traders} (commerçants) & \eng{to collect} (ramasser) & \eng{basket} (panier) \\
\eng{workers} (travailleurs) & \eng{to carry} (porter) & \eng{water can} (arrosoir / bidon) \\
\eng{volunteers} (bénévoles) & \eng{to plant} (planter) & \eng{shovel} (pelle) \\
\eng{youths} (jeunes) & \eng{to dig} (creuser) & \eng{pickaxe} (pioche) \\
\hline
\end{tabularx}
\end{center}
\vspace{0.3cm}
\begin{center}
\begin{tabularx}{\linewidth}{|l|X|l|}
\hline
\rowcolor{popgreenL} \textbf{Lieux} & \textbf{Temps / Fréquence} & \textbf{Résultats / Bénéfices} \\
\hline
\eng{market square} (place du marché) & \eng{early morning} (grand matin) & \eng{clean environment} (environnement propre) \\
\eng{health centre} (centre de santé) & \eng{break time} (pause) & \eng{healthy village} (village sain) \\
\eng{primary school} (école primaire) & \eng{noon} (midi) & \eng{new well} (nouveau puits) \\
\eng{road to the market} (route du marché) & \eng{monthly} (mensuel) & \eng{unity} (unité) \\
\eng{community radio} (radio communautaire) & \eng{annually} (annuel) & \eng{solidarity} (solidarité) \\
\hline
\end{tabularx}
\end{center}
\end{definition}

\begin{savaistu}[Le \eng{“Ikpwa”} et le travail communautaire au Cameroun]
Au Cameroun, le travail communautaire porte des noms selon les régions : \eng{Ikpwa} chez les Bassa, \eng{Njangi} ou \eng{Meetings} dans les Grassfields, \eng{Appel au travail} dans les zones urbaines. C'est une tradition ancestrale de solidarité (\eng{mutual aid}) où chaque ménage fournit une main-d'œuvre pour l'intérêt collectif : entretien des routes, construction d'écoles, de centres de santé, nettoyage des marchés. Le chef de quartier (\eng{quarter head}) ou le chef traditionnel mobilise la population. La participation est souvent obligatoire pour les résidents ; les absents s'acquittent d'une amende (\eng{fine}). C'est un vecteur puissant de cohésion sociale et de développement local.
\end{savaistu}

\courssub{Le relevé d'informations : compléter une fiche d'écoute}

La fiche d'écoute (\eng{listening grid}) est l'exercice le plus proche de la vraie vie : on écoute pour \textbf{noter} des informations précises, comme un journaliste ou un secrétaire de réunion. Chaque colonne correspond à une catégorie d'information ; une ou deux écoutes suffisent si l'on sait quoi chercher.

\begin{popfigure}
\begin{tikzpicture}
\draw[fill=popblueL, draw=popblue] (0,0) rectangle (11.5,1);
\node[align=center, font=\bfseries\small, text=popdark] at (5.75,0.5) {Listening Grid — Communal Work in Mbankomo};
\draw[fill=poporangeL, draw=poporange] (0,-1) rectangle (2.3,0);
\node[align=center, font=\bfseries\small, text=popdark] at (1.15,-0.5) {When?};
\draw[fill=popgreenL, draw=popgreen] (2.3,-1) rectangle (4.6,0);
\node[align=center, font=\bfseries\small, text=popdark] at (3.45,-0.5) {Where?};
\draw[fill=poppurpleL, draw=poppurple] (4.6,-1) rectangle (6.9,0);
\node[align=center, font=\bfseries\small, text=popdark] at (5.75,-0.5) {Who?};
\draw[fill=poppinkL, draw=poppink] (6.9,-1) rectangle (9.2,0);
\node[align=center, font=\bfseries\small, text=popdark] at (8.05,-0.5) {Tasks / Tools};
\draw[fill=popgoldL, draw=popgold] (9.2,-1) rectangle (11.5,0);
\node[align=center, font=\bfseries\small, text=popdark] at (10.35,-0.5) {Future plan};
\draw[draw=popline] (0,-3) rectangle (2.3,-1);
\draw[draw=popline] (2.3,-3) rectangle (4.6,-1);
\draw[draw=popline] (4.6,-3) rectangle (6.9,-1);
\draw[draw=popline] (6.9,-3) rectangle (9.2,-1);
\draw[draw=popline] (9.2,-3) rectangle (11.5,-1);
\node[align=center, font=\small, text=popmuted] at (1.15,-2) {date, time};
\node[align=center, font=\small, text=popmuted] at (3.45,-2) {place};
\node[align=center, font=\small, text=popmuted] at (5.75,-2) {number, leader};
\node[align=center, font=\small, text=popmuted] at (8.05,-2) {activities, tools};
\node[align=center, font=\small, text=popmuted] at (10.35,-2) {next project};
\end{tikzpicture}
\legende{Fiche d'écoute à compléter pendant l'écoute : une case = une information essentielle.}
\end{popfigure}

Voici la fiche complétée à titre de modèle de correction :

\begin{center}
\begin{tabularx}{\linewidth}{|l|X|}
\hline
\rowcolor{popblueL} Rubrique & Information relevée \\
\hline
\eng{When?} & \eng{Last Saturday, from 7 a.m. to noon.} \\
\hline
\eng{Where?} & \eng{Mbankomo (road to the market, market square, health centre).} \\
\hline
\eng{Who?} & \eng{More than 200 villagers; organised by the quarter head, Mr. Etoundi.} \\
\hline
\eng{Tasks?} & \eng{Clearing the road, sweeping the square, collecting rubbish, planting flowers.} \\
\hline
\eng{Tools?} & \eng{Machetes, hoes, baskets.} \\
\hline
\eng{Future plan?} & \eng{Next month: digging a new well near the primary school.} \\
\hline
\end{tabularx}
\end{center}

\courssub{La correction collective : reformuler en phrases complètes}

Après le travail individuel, le professeur corrige au tableau. Votre rôle n'est pas seulement de vérifier vos réponses : c'est d'apprendre à les \textbf{reformuler en phrases complètes}, compétence indispensable pour l'expression orale et écrite. Une note de fiche devient une phrase :

\begin{exemplebox}[De la note à la phrase]
Note de la fiche : \eng{When — last Saturday, 7 a.m. to noon.}

Phrase complète : \eng{The communal work took place last Saturday, from seven o'clock in the morning to noon.} « Les travaux communautaires ont eu lieu samedi dernier, de sept heures du matin à midi. »

Note : \eng{Future plan — new well, primary school.}

Phrase complète : \eng{Next month, the community will dig a new well near the primary school.} « Le mois prochain, la communauté creusera un nouveau puits près de l'école primaire. »
\end{exemplebox}

\courssub{Production guidée : réemploi du lexique}

Maintenant, à vous de parler de votre propre environnement en réemployant le vocabulaire et les structures vus.

\begin{experience}[Jeu de rôle : Interview radio — “Notre quartier en action”]
\textbf{Consigne :} En binôme. L'élève A est journaliste à la radio communautaire ; l'élève B est le chef de quartier (\eng{quarter head}) ou un habitant actif.
\begin{enumerate}
\item Le journaliste annonce le prochain travail communautaire (date, lieu, objectif).
\item Le chef de quartier explique les tâches prévues, les outils nécessaires, qui fait quoi.
\item Le journaliste demande pourquoi ce travail est important (bénéfices).
\item Le chef de quartier lance un appel à la participation.
\end{enumerate}
\textbf{Expressions utiles :} \eng{“We are organising… next Saturday.”} — \eng{“We need volunteers to…”} — \eng{“Bring your machetes and hoes.”} — \eng{“It will help us to…”} — \eng{“Everyone is welcome.”}
\end{experience}

\begin{exercice}[Expression écrite : Compte-rendu de travaux]
Rédigez un court paragraphe (60–80 mots) pour le journal de l'école ou le panneau d'affichage du quartier, relatant les derniers travaux communautaires chez vous (ou imaginez-les à partir de Mbankomo). Utilisez le \textbf{prétérit} pour les faits passés et le \textbf{futur} (\eng{will}) pour le projet à venir. Structurez votre texte :
\begin{enumerate}
\item Introduction : quand, qui a organisé, combien de participants.
\item Déroulement : tâches par groupes (hommes, femmes, jeunes), outils.
\item Moment de convivialité (pause, repas).
\item Projet futur annoncé + raison.
\end{enumerate}
\end{exercice}

\begin{exercice}[À vous de jouer — Compréhension]
Répondez aux questions suivantes sur le script de Mbankomo. Pour les Vrai/Faux, justifiez par une citation.
\begin{enumerate}
\item \eng{True or False: The women cleared the road with machetes.}
\item \eng{True or False: The traders offered food to the workers.}
\item \eng{Who announced the work on the radio? a) the mayor \quad b) the traders \quad c) the quarter head}
\item \eng{What did the girls plant around the health centre?}
\item \eng{Why will the community dig a new well?}
\end{enumerate}
\end{exercice}

\begin{corrige}[Exercice « À vous de jouer »]
\begin{enumerate}
\item \eng{False. The text says: “The men cleared the road to the market with machetes and hoes.”} (Ce sont les hommes, pas les femmes.)
\item \eng{True. The text says: “The traders of the market offered food and palm wine to the workers.”}
\item Réponse \textbf{c) the quarter head} : \eng{The quarter head, Mr. Etoundi, had announced the work on the community radio.}
\item \eng{Flowers.} (Phrase complète : \eng{The girls planted flowers around the health centre.})
\item \eng{Because the old well dries up during the dry season.}
\end{enumerate}
\end{corrige}

\begin{aretenir}[Les quatre réflexes du bon auditeur]
\begin{itemize}
\item \eng{True / False} : toujours justifier par une citation du texte, sinon vous perdez des points.
\item QCM : méfiez-vous des mots du texte placés en leurres ; vérifiez \emph{qui fait quoi}.
\item Réponses courtes : répondez directement, sans recopier la question.
\item Fiche d'écoute : une case = une information ; reformulez ensuite en phrases complètes.
\end{itemize}
\end{aretenir}

Ces stratégies de compréhension et ce vocabulaire vous serviront dans toutes les épreuves d'écoute de l'année. Entraînez-vous à écouter des actualités locales (CRTV, radios communautaires) pour repérer les mêmes structures : \eng{who, what, when, where, why, how many}.

\coursec{Lexique thématique du travail communautaire}

Dans la section précédente, vous avez répondu à des questions de compréhension sur le script d'écoute : vous avez repéré le thème, les mots-clés, les chiffres et les noms propres. Vous avez certainement remarqué que certaines expressions revenaient souvent : \eng{communal work}, \eng{volunteers}, \eng{spades}, \eng{to sweep the streets}. C'est exactement ce que nous allons étudier maintenant de façon systématique : le \cle{vocabulaire} du travail communautaire. Un bon lexique, bien organisé en familles de mots, est la clé pour comprendre un enregistrement sans paniquer et pour réemployer les mots dans vos propres phrases, à l'écrit comme à l'oral.

\courssub{Observer avant d'apprendre : le lexique en contexte}

Avant de mémoriser des listes, observons comment ces mots vivent dans un texte authentique. Lisez le texte suivant et soulignez mentalement tous les mots qui désignent des outils, des tâches, des personnes et des valeurs.

\begin{textebox}[Keeping Molyko clean — a report from Buea]
Every last Saturday of the month, the people of Molyko, a lively quarter of Buea, take part in communal work. At six o'clock in the morning, the quarter head beats the traditional drum, and villagers of all ages come out with their tools. Some carry wheelbarrows and spades; others bring hoes, cutlasses, brooms, rakes and buckets.

The tasks are shared fairly. The young men clear the gutters so that rain water can flow freely during the rainy season. The women sweep the streets in front of the market, while a group of volunteers repairs the small wooden bridge near the river. Last month, the community also dug a well behind the health centre, and the pupils of Government School Molyko planted trees along the main road. A few farmers offered to weed the grass around the primary school.

Nobody is paid for this work. A local trader provides drinking water and bread for everybody, and the council sometimes lends tools such as wheelbarrows and rakes. The traditional ruler himself often joins the workers, which shows that communal work concerns everybody, from the youngest pupil to the most respected elder.

“Communal work teaches us unity and solidarity,” says Mrs Eposi, a mother of four. “When we work together, we keep our village clean, and everybody benefits from it. A clean quarter means fewer mosquitoes, less disease and more pride.”

At eleven o'clock, the tools are returned, the streets are shining, and neighbours share food and palm wine. Communal work is not only about cleaning: it is about building a family larger than every household.
\end{textebox}

\textbf{Glossaire (mots difficiles)}\par
\begin{itemize}
\item \eng{quarter} (n.) : quartier ; \eng{quarter head} : chef de quartier
\item \eng{gutter} (n.) : caniveau ; \eng{to clear} : déboucher, nettoyer
\item \eng{well} (n.) : puits ; \eng{to dig} : creuser
\item \eng{to weed} (v.) : désherber ; \eng{elder} (n.) : aîné, ancien
\item \eng{to lend} (v.) : prêter (à quelqu'un) ; \eng{to benefit from} : profiter de
\item \eng{household} (n.) : ménage, foyer ; \eng{palm wine} : vin de palme
\end{itemize}

\textbf{Questions rapides (répondez en phrases courtes)}\par
\begin{enumerate}
\item Name four tools mentioned in the text.
\item Who provides food and drinks for the workers?
\item What did the pupils of Government School Molyko do?
\item According to Mrs Eposi, what are the benefits of communal work?
\end{enumerate}

\begin{corrige}[Questions rapides]
\begin{enumerate}
\item \eng{Wheelbarrows, spades, hoes and brooms are mentioned.} (toute combinaison de quatre outils du texte est acceptée : \eng{cutlasses, rakes, buckets})
\item \eng{A local trader provides drinking water and bread for everybody.}
\item \eng{They planted trees along the main road.}
\item \eng{According to Mrs Eposi, communal work teaches unity and solidarity, keeps the village clean, and means fewer mosquitoes, less disease and more pride.}
\end{enumerate}
\end{corrige}

\courssub{Champ lexical 1 : les outils (tools)}

\begin{definition}[Tools — les outils]
Un \cle{outil} se dit \eng{a tool} (prononcé « toul »). Le mot est \textbf{comptable} : \eng{a tool, two tools}. L'ensemble du matériel se dit \eng{tools} ou \eng{equipment} (ce dernier est \textbf{incomptable} : jamais de \emph{s}, jamais d'article \eng{a} devant).
\end{definition}

\begin{center}
\begin{tabularx}{\linewidth}{|X|X|X|}
\hline
\rowcolor{popblueL} English & Français & Exemple \\
\hline
a wheelbarrow & une brouette & \eng{The men pushed two wheelbarrows full of sand.} \\
\hline
a spade & une pelle & \eng{She used a spade to dig the hole.} \\
\hline
a hoe & une houe & \eng{Farmers weed the farm with a hoe.} \\
\hline
a cutlass & une machette & \eng{He cut the tall grass with his cutlass.} \\
\hline
a broom & un balai & \eng{Take a broom and sweep the veranda.} \\
\hline
a rake & un râteau & \eng{They gathered the dead leaves with a rake.} \\
\hline
a bucket & un seau & \eng{She carried a bucket of water on her head.} \\
\hline
\end{tabularx}
\end{center}

\begin{attention}[Prononciation et orthographe]
\eng{wheelbarrow} se prononce « ouil-bé-ro » : ne dites pas « ouil-bar-rou ». Attention aussi au pluriel : \eng{two spades}, \eng{three buckets} — le \emph{s} du pluriel est obligatoire en anglais, même si en français on entend parfois « deux pelle » à l'oral. Enfin, \eng{cutlass} est le mot courant en Afrique anglophone et dans les Caraïbes ; \eng{machete} existe aussi mais est d'origine espagnole.
\end{attention}

\courssub{Champ lexical 2 : les tâches (tasks)}

Les tâches s'expriment avec des \cle{verbes d'action} suivis d'un complément. Apprenez chaque verbe \textbf{avec} son complément habituel : c'est ce qu'on appelle une \cle{collocation}.

\begin{center}
\begin{tabularx}{\linewidth}{|X|X|X|}
\hline
\rowcolor{popblueL} English & Français & Exemple traduit \\
\hline
to clear the gutters & déboucher les caniveaux & \eng{They cleared the gutters before the rains.} « Ils ont débouché les caniveaux avant les pluies. » \\
\hline
to sweep the streets & balayer les rues & \eng{We swept the streets on Saturday.} « Nous avons balayé les rues samedi. » \\
\hline
to repair a bridge & réparer un pont & \eng{The volunteers repaired the bridge.} « Les volontaires ont réparé le pont. » \\
\hline
to dig a well & creuser un puits & \eng{The village dug a well last year.} « Le village a creusé un puits l'an dernier. » \\
\hline
to plant trees & planter des arbres & \eng{Pupils planted trees along the road.} « Les élèves ont planté des arbres le long de la route. » \\
\hline
to weed & désherber & \eng{She weeded the school garden.} « Elle a désherbé le jardin de l'école. » \\
\hline
\end{tabularx}
\end{center}

\begin{propriete}[Verbes irréguliers à connaître]
Trois verbes de ce champ lexical sont irréguliers. Apprenez leurs trois formes :
\begin{center}
\begin{tabular}{|l|l|l|l|}
\hline
\rowcolor{popblueL} Base verbale & Prétérit & Participe passé & Français \\
\hline
sweep & swept & swept & balayer \\
\hline
dig & dug & dug & creuser \\
\hline
take (part in) & took & taken & prendre (part à) \\
\hline
\end{tabular}
\end{center}
Les autres sont réguliers : \eng{repair — repaired}, \eng{plant — planted}, \eng{weed — weeded}, \eng{clear — cleared}. Ajoutez aussi \eng{lend — lent — lent} (prêter) et \eng{bring — brought — brought} (apporter), très fréquents dans ce thème.
\end{propriete}

\begin{attention}[Piège de traduction : « balayer »]
\eng{to sweep} signifie « balayer » et non « essuyer ». « Essuyer » se dit \eng{to wipe} (une table) ou \eng{to mop} (le sol avec une serpillière). Ne dites donc jamais \eng{sweep the floor with water} pour « laver le sol » : dites \eng{mop the floor}.
\end{attention}

\courssub{Champ lexical 3 : les personnes (people)}

\begin{definition}[Volunteer — volontaire]
A \cle{volunteer} is \emph{a person who works willingly without being paid} — une personne qui travaille volontairement sans être payée : un volontaire, un bénévole. Prononciation : « vo-len-tière » (accent sur la dernière syllabe).
\end{definition}

\begin{center}
\begin{tabularx}{\linewidth}{|X|X|X|}
\hline
\rowcolor{popblueL} English & Français & Exemple \\
\hline
the quarter head & le chef de quartier & \eng{The quarter head announced the date of the work.} \\
\hline
the traditional ruler & le chef traditionnel & \eng{The traditional ruler joined the workers.} \\
\hline
a villager & un(e) villageois(e) & \eng{Every villager brought a tool.} \\
\hline
a volunteer & un volontaire, un bénévole & \eng{Twenty volunteers repaired the bridge.} \\
\hline
a trader & un(e) commerçant(e) & \eng{A trader provided bread and water.} \\
\hline
the council & le conseil municipal, la municipalité & \eng{The council lent wheelbarrows.} \\
\hline
\end{tabularx}
\end{center}

\begin{attention}[Faux ami : « communal » et « council »]
Attention à deux pièges classiques :
\begin{itemize}
\item \eng{communal} ne veut pas dire « communal » au sens français de « qui appartient à la mairie ». \eng{communal work} signifie \textbf{travail collectif} organisé par la communauté du quartier ou du village, pas forcément par la mairie.
\item \eng{the council} désigne la \textbf{municipalité, le conseil municipal} (l'institution). Ne le confondez pas avec \eng{counsel} (conseil, avis) ni avec le verbe français « conseiller » (\eng{to advise}).
\end{itemize}
Ainsi : \eng{The council provided spades and cutlasses.} « La commune a fourni des pelles et des machettes. » — mais \eng{The villagers organised communal work.} « Les villageois ont organisé un travail collectif. »
\end{attention}

\courssub{Champ lexical 4 : organisation et valeurs (values)}

\begin{center}
\begin{tabularx}{\linewidth}{|X|X|X|}
\hline
\rowcolor{popblueL} English & Français & Remarque \\
\hline
to take part in & participer à & suivi d'un nom : \eng{take part in communal work} \\
\hline
to join & rejoindre, se joindre à & \textbf{sans} préposition : \eng{join the workers} \\
\hline
to provide & fournir & \eng{provide something} ou \eng{provide somebody with something} \\
\hline
unity & l'unité & nom incomptable \\
\hline
solidarity & la solidarité & nom incomptable \\
\hline
a benefit & un bénéfice, un avantage & comptable ; verbe : \eng{to benefit from} \\
\hline
\end{tabularx}
\end{center}

\begin{attention}[« Participer » : to take part in, pas « to participate in » à l'oral]
Le verbe \eng{to participate in} existe, mais il est soutenu. À l'oral et dans les textes courants, on préfère \eng{to take part in}. Deuxième piège : \eng{to join} ne prend \textbf{jamais} la préposition \eng{in} : on dit \eng{I joined the volunteers} (« J'ai rejoint les volontaires »), jamais \eng{I joined in the volunteers}. Enfin, \eng{to assist} signifie « aider », pas « assister à » : pour « assister à une réunion », dites \eng{to attend a meeting}.
\end{attention}

\begin{exemplebox}[Les collocations à retenir par cœur]
\begin{itemize}
\item \eng{to take part in communal work} : participer au travail communautaire
\item \eng{to carry out a project} : mener à bien un projet — \eng{The youth carried out a project to clean the market.} « Les jeunes ont mené un projet de nettoyage du marché. »
\item \eng{to work together} : travailler ensemble — \eng{We work together every last Saturday.} « Nous travaillons ensemble chaque dernier samedi du mois. »
\item \eng{to keep the village clean} : garder le village propre — \eng{We must keep our village clean.} « Nous devons garder notre village propre. »
\end{itemize}
\end{exemplebox}

\courssub{Carte mentale : organiser le lexique pour mieux le mémoriser}

Une carte mentale regroupe les mots par familles. Retenez quatre mots par branche : c'est une excellente stratégie de mémorisation avant un test d'écoute.

\begin{popfigure}
\begin{tikzpicture}[mindmap, grow cyclic, every node/.style={concept}, text=white,
root concept/.append style={concept color=popblue, font=\small\bfseries},
level 1/.append style={level distance=3.1cm, sibling angle=90},
level 2/.append style={level distance=2.1cm, sibling angle=45, concept color=popmuted, font=\tiny}]
\node[root concept]{COMMUNAL WORK}
  child[concept color=poporange]{ node{TOOLS}
    child{ node{wheelbarrow} }
    child{ node{spade} }
    child{ node{cutlass} }
    child{ node{broom} }
  }
  child[concept color=popgreen]{ node{TASKS}
    child{ node{clear gutters} }
    child{ node{sweep streets} }
    child{ node{dig a well} }
    child{ node{plant trees} }
  }
  child[concept color=poppurple]{ node{PEOPLE}
    child{ node{quarter head} }
    child{ node{villager} }
    child{ node{volunteer} }
    child{ node{council} }
  }
  child[concept color=poppink]{ node{VALUES}
    child{ node{unity} }
    child{ node{solidarity} }
    child{ node{benefit} }
    child{ node{take part in} }
  };
\end{tikzpicture}
\legende{Carte mentale du lexique du travail communautaire : quatre familles de mots à mémoriser.}
\end{popfigure}

\courssub{Les familles de mots : un mot, plusieurs formes}

En anglais comme en français, un même radical donne un nom, un adjectif, un adverbe. Connaître ces familles multiplie votre vocabulaire sans effort supplémentaire.

\begin{center}
\begin{tabular}{|l|l|l|}
\hline
\rowcolor{popblueL} Nom & Adjectif / verbe & Adverbe / nom d'action \\
\hline
a volunteer (volontaire) & voluntary (bénévole) & voluntarily (bénévolement) \\
\hline
a participant (participant) & to participate (participer) & participation (participation) \\
\hline
a benefit (bénéfice) & beneficial (bénéfique) & to benefit from (profiter de) \\
\hline
\end{tabular}
\end{center}

\begin{exemplebox}[La famille en action]
\begin{itemize}
\item \eng{Fifty volunteers came on Saturday.} « Cinquante volontaires sont venus samedi. » (nom)
\item \eng{Communal work is voluntary: nobody forces you.} « Le travail communautaire est bénévole : personne ne vous oblige. » (adjectif)
\item \eng{She voluntarily offered her wheelbarrow.} « Elle a offert sa brouette de son plein gré. » (adverbe)
\item \eng{Everybody's participation is important.} « La participation de chacun est importante. » (nom d'action)
\end{itemize}
\end{exemplebox}

\begin{attention}[Attention à la prononciation de « volunteer »]
Le mot \eng{volunteer} porte l'accent sur la \textbf{dernière} syllabe : « vo-len-TIERE ». En revanche, l'adjectif \eng{voluntary} porte l'accent sur la \textbf{première} syllabe : « VO-len-teu-ri ». Un mauvais accent peut empêcher un anglophone de vous comprendre.
\end{attention}

\begin{exoresolu}[Complétez avec le mot qui convient]
Énoncé — Complete each sentence with a word from the lesson: \emph{spades, volunteers, take part in, gutters, keep, council, solidarity}.
\begin{enumerate}
\item The \ldots\ldots lent us ten wheelbarrows for the clean-up day.
\item We use \ldots\ldots to dig holes for the new trees.
\item Every Saturday, the villagers \ldots\ldots communal work.
\item Before the rainy season, we must clear the \ldots\ldots.
\item Communal work teaches us unity and \ldots\ldots.
\item We must \ldots\ldots our quarter clean.
\item Thirty \ldots\ldots repaired the bridge without being paid.
\end{enumerate}
\tcblower
Solution détaillée :
\begin{enumerate}
\item \textbf{council} — c'est la municipalité qui prête le matériel (\eng{lend — lent — lent}).
\item \textbf{spades} — on creuse des trous avec des pelles ; le pluriel est naturel ici car on parle des pelles en général.
\item \textbf{take part in} — collocation : \eng{take part in communal work} ; présent simple avec un sujet pluriel.
\item \textbf{gutters} — \eng{clear the gutters} : déboucher les caniveaux.
\item \textbf{solidarity} — nom de valeur, parallèle à \eng{unity}.
\item \textbf{keep} — collocation : \eng{keep our quarter clean} ; après \eng{must}, la base verbale.
\item \textbf{volunteers} — des personnes qui travaillent sans être payées ; pluriel obligatoire après \eng{thirty}.
\end{enumerate}
\end{exoresolu}

\begin{exercice}[À vous de jouer : réemploi du lexique]
\begin{enumerate}
\item Traduisez en anglais : « Les villageois ont creusé un puits et balayé les rues samedi dernier. »
\item Complétez : \eng{The pupils \ldots\ldots (join) the workers without being paid.} (mettez le verbe au prétérit)
\item Trouvez l'intrus : \eng{wheelbarrow — spade — solidarity — rake}. Justifiez en une phrase en anglais.
\item Rédigez deux phrases sur le travail communautaire dans votre quartier, avec \eng{take part in} et \eng{keep \ldots\ clean}.
\end{enumerate}
\end{exercice}

\begin{corrige}[Exercice : réemploi du lexique]
\begin{enumerate}
\item \eng{The villagers dug a well and swept the streets last Saturday.} (prétérits irréguliers : \eng{dig — dug}, \eng{sweep — swept})
\item \eng{The pupils joined the workers without being paid.} (\eng{join} est régulier : \eng{joined} ; sans préposition.)
\item L'intrus est \eng{solidarity}. \eng{Solidarity is a value, but a wheelbarrow, a spade and a rake are tools.}
\item Modèle : \eng{Every month, I take part in communal work in my quarter. We sweep the streets and keep our neighbourhood clean.}
\end{enumerate}
\end{corrige}

\begin{savaistu}[Le « nnoboa » au Ghana]
Au Ghana, le travail communautaire est une tradition ancestrale appelée \eng{nnoboa} en langue akan : l'entraide mutuelle entre voisins. Quand une famille doit construire une maison, labourer un champ ou organiser des funérailles, tout le village vient l'aider gratuitement, et la famille aidée fera de même pour les autres. On retrouve la même idée au Cameroun anglophone avec les \eng{community work days} et au Cameroun francophone avec le « travail communautaire » du village. Partout en Afrique de l'Ouest, travailler ensemble est une valeur aussi importante que l'école ou le marché !
\end{savaistu}

\begin{aretenir}[L'essentiel du lexique]
\begin{itemize}
\item Quatre familles de mots : \textbf{tools} (\eng{wheelbarrow, spade, hoe, cutlass, broom, rake, bucket}), \textbf{tasks} (\eng{clear the gutters, sweep the streets, repair a bridge, dig a well, plant trees, weed}), \textbf{people} (\eng{quarter head, traditional ruler, villager, volunteer, trader, council}), \textbf{values} (\eng{unity, solidarity, benefit}).
\item Collocations clés : \eng{take part in communal work}, \eng{carry out a project}, \eng{work together}, \eng{keep the village clean}.
\item Familles de mots : \eng{volunteer / voluntary / voluntarily} ; \eng{participate / participation / participant}.
\item Pièges : \eng{communal} = collectif (pas « de la mairie ») ; \eng{join} sans préposition ; \eng{sweep} = balayer (pas essuyer) ; \eng{sweep — swept — swept}, \eng{dig — dug — dug}, \eng{lend — lent — lent}.
\end{itemize}
\end{aretenir}

Maintenant que ce lexique est bien organisé dans votre tête et dans votre cahier, vous êtes prêt pour la suite : la méthode d'écoute, où vous apprendrez à repérer ces mots-clés rapidement dans un enregistrement, puis la production guidée, où vous les réemploierez dans vos propres phrases.

\coursec{Méthode d'écoute et exemple traité}

Dans la section précédente, nous avons construit le vocabulaire du travail communautaire : \eng{to sweep the streets}, \eng{to clear the gutters}, \eng{to fetch water}, \eng{to weed the school compound}... Mais connaître les mots ne suffit pas : encore faut-il les \emph{reconnaître à l'oreille}, au milieu d'un enregistrement lu à vitesse normale. Cette section vous donne une méthode complète, en quatre étapes, pour réussir n'importe quelle épreuve de \cle{listening comprehension}, puis un exemple traité pas à pas sur le script de Bafut.

\courssub{Vue d'ensemble : les quatre étapes de l'écoute efficace}

Une écoute réussie ne commence pas quand l'enregistrement démarre : elle commence \emph{avant}, avec la lecture des questions, et se termine \emph{après}, avec la vérification des réponses. Voici la carte complète de la méthode.

\begin{popfigure}
\begin{tikzpicture}[
  box/.style={draw=popblue, thick, fill=popblueL, rounded corners=3pt, align=center, text width=3.1cm, minimum height=1.5cm, font=\small},
  arr/.style={-{Stealth[length=3mm]}, very thick, poporange}
]
\node[box, align=center] (b1) at (0,0){\textbf{BEFORE}\\ Lire les questions,\\ souligner les mots-clés,\\ \emph{predict}};
\node[box, right=1.1cm of b1, align=center] (b2){\textbf{LISTEN 1}\\ Capter le sens global :\\ \emph{who, what, where}\\ (\emph{gist})};
\node[box, right=1.1cm of b2, align=center] (b3){\textbf{LISTEN 2}\\ Noter les détails :\\ chiffres, heures, noms,\\ tâches};
\node[box, right=1.1cm of b3, align=center] (b4){\textbf{AFTER}\\ Rédiger en phrases\\ complètes, vérifier\\ l'orthographe};
\draw[arr] (b1) -- (b2);
\draw[arr] (b2) -- (b3);
\draw[arr] (b3) -- (b4);
\end{tikzpicture}
\legende{Les quatre étapes d'une écoute efficace : prédire, capter le sens global, noter les détails, rédiger et vérifier.}
\end{popfigure}

\courssub{Étape 1 — Avant l'écoute : prédire}

\begin{propriete}[Lire les questions avant d'écouter]
Avant la première écoute, lisez \textbf{toutes} les questions. Pour chacune :
\begin{enumerate}
\item \textbf{soulignez les mots-clés} (le mot interrogatif et le sujet) : \eng{\textbf{What} did the \textbf{women} do?} ;
\item \textbf{identifiez le type d'information attendu} : une action (\eng{what... do}), un lieu (\eng{where}), une heure (\eng{when}, \eng{what time}), un nombre (\eng{how many}), une personne (\eng{who}) ;
\item \textbf{prédisez les réponses possibles} grâce au thème : si le sujet est le travail communautaire, la réponse à \eng{What did the women do?} sera probablement une tâche du type \eng{sweep}, \eng{clean}, \eng{cook}, \eng{fetch water}.
\end{enumerate}
\end{propriete}

La prédiction n'est pas de la devinette : c'est une \emph{anticipation} qui prépare votre oreille. Si vous attendez un nombre, vous serez attentif à tous les chiffres ; si vous attendez un lieu, vous guetterez les noms de villages et de quartiers.

\begin{exemplebox}[Prédire à partir d'une question]
Question : \eng{How many villagers took part in the work?}

Mots-clés : \eng{how many} + \eng{villagers}. Type attendu : un \textbf{nombre}. Prédictions possibles : \eng{fifty}, \eng{about a hundred}, \eng{more than two hundred}. Pendant l'écoute, toute quantité prononcée devient un signal : \eng{over two hundred people...} — voilà la réponse.
\end{exemplebox}

\courssub{Étape 2 — Première écoute : capter le sens global}

\begin{methode}[Première écoute : le « gist »]
\begin{enumerate}
\item Écoutez \textbf{sans rien écrire} (ou presque) : ne perdez pas le fil en prenant des notes trop tôt.
\item Répondez mentalement à trois questions : \eng{Who?} (qui parle ? de qui parle-t-on ?), \eng{What?} (de quoi s'agit-il ?), \eng{Where?} (où cela se passe-t-il ?).
\item Formulez le thème en une phrase simple : \eng{A report about communal work in Bafut.}
\item Notez éventuellement les noms propres dès leur première occurrence (voir la boîte Attention plus bas).
\end{enumerate}
\end{methode}

\begin{attention}[Ne cherchez pas à comprendre chaque mot]
C'est l'erreur numéro un des élèves francophones : paniquer dès qu'un mot est inconnu et « décrocher » pendant dix secondes. Un enregistrement, même au Bac, contient toujours des mots que vous ne connaissez pas. \textbf{Concentrez-vous sur les mots porteurs de sens} : les noms, les verbes, les chiffres, les noms propres. Le reste (articles, petits mots de liaison) peut être ignoré sans dommage. Si vous entendez \eng{the women \textbf{swept} the \textbf{streets}}, vous avez tout compris, même si un adjectif vous a échappé.
\end{attention}

\courssub{Étape 3 — Deuxième écoute : noter les détails}

La deuxième écoute est une écoute \emph{ciblée} : vous savez déjà de quoi parle le texte, vous cherchez maintenant les informations précises demandées par les questions.

\begin{propriete}[Que faut-il noter à la deuxième écoute ?]
\begin{itemize}
\item les \textbf{chiffres et quantités} : \eng{over 200 villagers}, \eng{three trucks of sand} ;
\item les \textbf{heures et durées} : \eng{from 7 a.m. to noon} ;
\item les \textbf{noms propres} : personnes, villages, écoles, associations ;
\item les \textbf{tâches et actions} : qui a fait quoi (\eng{the men cleared the gutters}) ;
\item les \textbf{dates et jours} : \eng{last Saturday}, \eng{every first Saturday of the month}.
\end{itemize}
\end{propriete}

\courssub{La technique de prise de notes : abréviations, flèches et symboles}

Vous n'avez pas le temps d'écrire des mots entiers pendant l'écoute. Utilisez un système personnel d'\cle{abréviations} et de symboles, que vous « traduirez » en phrases complètes après l'écoute.

\begin{center}
\begin{tabularx}{\linewidth}{|l|X|X|}\hline
\rowcolor{popblueL} \textbf{Note rapide} & \textbf{Signification} & \textbf{Phrase rédigée ensuite} \\ \hline
Sat. & Saturday & \eng{The work took place on Saturday.} \\ \hline
vill. & villagers & \eng{The villagers cleaned the market square.} \\ \hline
200+ & more than two hundred & \eng{More than two hundred people took part.} \\ \hline
7$\rightarrow$12 & from 7 a.m. to noon & \eng{They worked from 7 a.m. to noon.} \\ \hline
W $\rightarrow$ sweep & the women swept & \eng{The women swept the streets.} \\ \hline
M $\rightarrow$ gutters & the men cleared the gutters & \eng{The men cleared the gutters.} \\ \hline
bec. / so & because / so, therefore & \eng{They worked together because they wanted a clean village.} \\ \hline
1st & first & \eng{It happens every first Saturday of the month.} \\ \hline
\end{tabularx}
\end{center}

\begin{exemplebox}[De la note à la phrase complète]
Pendant l'écoute, vous entendez : \eng{The exercise started at seven in the morning and ended at noon.}

Vous notez : \textbf{7$\rightarrow$12}.

Après l'écoute, vous rédigez : \eng{They worked from 7 a.m. to noon.} « Ils ont travaillé de 7 heures du matin à midi. »
\end{exemplebox}

\courssub{Étape 4 — Après l'écoute : rédiger et vérifier}

\begin{propriete}[Rédiger les réponses]
\begin{enumerate}
\item Répondez par des \textbf{phrases complètes} (sujet + verbe), sauf si la consigne autorise les réponses courtes. À la question \eng{What did the women do?}, écrivez \eng{They swept the streets.} et non simplement \eng{swept the streets}.
\item \textbf{Recyclez les mots de la question} dans votre réponse : cela garantit la bonne structure et le bon temps verbal.
\item Vérifiez l'\textbf{orthographe des noms propres} (\eng{Bafut}, \eng{Mr Ngum}, \eng{Mrs Abong}) et la \textbf{conjugaison} : si la question est au prétérit (\eng{did}), la réponse est au prétérit (\eng{swept}, pas \eng{sweep}).
\end{enumerate}
\end{propriete}

\begin{attention}[Les noms propres : notez-les dès la première occurrence]
Dans les enregistrements d'examen, les noms propres — \eng{Bafut}, \eng{Mr Ngum}, \eng{Mrs Abong} — sont souvent \textbf{épelés lettre par lettre} ou \textbf{répétés}. Notez-les immédiatement, phonétiquement si nécessaire (« ba-foute », « n-goum »), puis vérifiez l'orthographe si le locuteur épelle : \eng{N-G-U-M}. Un nom propre mal orthographié peut coûter le point de la question.
\end{attention}

\courssub{Les signaux phoniques à repérer}

En anglais parlé, l'information importante n'est pas cachée : elle est \emph{accentuée}.

\begin{propriete}[L'accent de phrase porte l'information]
\begin{itemize}
\item Les \textbf{mots accentués} (prononcés plus fort et plus longuement) portent le sens : \eng{The \textbf{women} \textbf{swept} the \textbf{streets}.} Les petits mots (\eng{the}, \eng{a}, \eng{of}) sont faibles et rapides.
\item Les \textbf{chiffres} sont souvent \textbf{répétés} ou suivis d'une \textbf{pause} : \eng{Over two hundred — yes, two hundred — villagers came.} Cette répétition est un cadeau : c'est presque toujours une réponse à une question.
\item Les mots comme \eng{first}, \eng{then}, \eng{finally}, \eng{but}, \eng{however} annoncent un changement ou une information nouvelle : redoublez d'attention juste après.
\item Le mot \eng{while} (« pendant que / tandis que ») introduit souvent un \textbf{contraste entre deux groupes} : \eng{The men cleared the gutters, \textbf{while} the women swept the streets.} C'est le signal qu'il faut noter deux actions différentes pour deux groupes différents.
\end{itemize}
\end{propriete}

\begin{savaistu}[L'enregistrement est lu deux fois : utilisez-les bien]
Au baccalauréat et dans la plupart des évaluations, l'enregistrement est généralement lu \textbf{deux fois}. La première écoute sert à comprendre le sens global — \textbf{jamais} à rédiger les réponses définitives. Beaucoup d'élèves écrivent trop vite une réponse à la première écoute, puis découvrent à la deuxième qu'elle était fausse ou incomplète... et n'ont plus le temps de corriger proprement. Première écoute : comprendre. Deuxième écoute : noter. Après : rédiger.
\end{savaistu}

\courssub{Exemple traité pas à pas}

Appliquons toute la méthode sur un extrait d'enregistrement typique.

\begin{textebox}[Listening script — Communal work in Bafut]
\eng{Good morning, listeners. Last Saturday, the people of Bafut carried out their monthly communal work. More than two hundred villagers took part in the exercise, which started at seven in the morning and ended at noon. The men cleared the gutters and repaired the bridge near the market, while the women swept the streets and collected the rubbish. The students of Government High School Bafut weeded the school compound. Mr Ngum, the quarter head, thanked everybody and announced that the next exercise will take place on the first Saturday of next month.}
\end{textebox}

Glossaire : \eng{to carry out} (v., réaliser, effectuer) ; \eng{gutter} (n., caniveau) ; \eng{rubbish} (n., ordures) ; \eng{to weed} (v., désherber) ; \eng{quarter head} (n., chef de quartier).

\begin{exoresolu}[Question : What did the women do?]
Répondez par une phrase complète, d'après l'enregistrement sur le travail communautaire de Bafut.
\tcblower
\textbf{Étape 1 — Avant l'écoute.} Mots-clés de la question : \eng{what} + \eng{women} + \eng{do}. Type d'information attendu : une ou plusieurs \textbf{actions} accomplies par les femmes. Prédiction (thème du travail communautaire) : probablement \eng{sweep}, \eng{clean}, \eng{cook} ou \eng{fetch water}.

\textbf{Étape 2 — Première écoute.} Sens global : un reportage sur le travail communautaire mensuel à Bafut. Je repère qu'on parle des hommes, des femmes et des élèves : la réponse viendra donc au moment où le locuteur mentionne \eng{the women}.

\textbf{Étape 3 — Deuxième écoute ciblée.} J'entends : \eng{...while the \textbf{women} \textbf{swept} the \textbf{streets} and \textbf{collected} the \textbf{rubbish}.} Je note rapidement : \textbf{W $\rightarrow$ sweep streets + collect rubbish}.

\textbf{Étape 4 — Rédaction.} Je rédige une phrase complète au prétérit, en recyclant les mots de la question :

\eng{The women swept the streets and collected the rubbish.}

« Les femmes ont balayé les rues et ramassé les ordures. »

Vérification : sujet + verbe au prétérit (\eng{swept}, \eng{collected}), orthographe correcte, information complète (deux actions, pas une seule).
\end{exoresolu}

\begin{attention}[Pièges fréquents des francophones]
\begin{itemize}
\item \textbf{Répondre au présent} alors que la question est au prétérit : \eng{The women sweep the streets} est faux ici ; l'événement est passé (\eng{last Saturday}), il faut \eng{swept} (prétérit irrégulier de \eng{sweep} : \eng{sweep -- swept -- swept}).
\item \textbf{Calquer le français} : ne dites pas \eng{the women made the streets clean} pour « les femmes ont nettoyé les rues » si le texte dit \eng{swept} ; reprenez le verbe de l'enregistrement.
\item \textbf{Oublier la moitié de la réponse} : quand le script donne deux actions reliées par \eng{and}, donnez les deux.
\item \textbf{Confondre les groupes} : les hommes ont réparé le pont, les femmes ont balayé, les élèves ont désherbé. Notez \eng{M}, \eng{W}, \eng{St} pour ne pas mélanger.
\item \textbf{Faux ami} : \eng{exercise} signifie ici « opération, activité » (le travail communautaire), pas « exercice scolaire ».
\end{itemize}
\end{attention}

\begin{exercice}[À vous de jouer]
Réécoutez mentalement le script de Bafut et répondez par des phrases complètes :
\begin{enumerate}
\item \eng{When did the communal work take place?}
\item \eng{How many villagers took part in the exercise?}
\item \eng{What did the men do?}
\item \eng{What did the students of Government High School Bafut do?}
\item \eng{When will the next exercise take place?}
\end{enumerate}
\end{exercice}

\begin{corrige}[Exercice — À vous de jouer]
\begin{enumerate}
\item \eng{It took place last Saturday.} (Mots-clés : \eng{when} $\rightarrow$ un moment ; signal entendu : \eng{last Saturday}. Prétérit irrégulier : \eng{take -- took -- taken}.)
\item \eng{More than two hundred villagers took part.} (Le chiffre est répété et suivi d'une pause : signal classique.)
\item \eng{The men cleared the gutters and repaired the bridge near the market.} (Deux actions : ne pas en oublier une.)
\item \eng{They weeded the school compound.} (\eng{weed} est un verbe régulier : \eng{weed -- weeded -- weeded}.)
\item \eng{It will take place on the first Saturday of next month.} (Attention au futur : la question est \eng{will}, la réponse garde \eng{will take place}.)
\end{enumerate}
\end{corrige}

\begin{aretenir}[L'essentiel de la méthode]
\begin{itemize}
\item \textbf{Avant} : lire les questions, souligner les mots-clés, prédire le type de réponse.
\item \textbf{Écoute 1} : sens global uniquement (\eng{who, what, where}) ; ne pas rédiger.
\item \textbf{Écoute 2} : noter les détails avec des abréviations et des flèches (\eng{7$\rightarrow$12}, \eng{W $\rightarrow$ sweep}).
\item \textbf{Après} : rédiger en phrases complètes, au bon temps, et vérifier les noms propres.
\item Les mots accentués portent l'information ; les chiffres répétés sont des réponses déguisées ; \eng{while} annonce un contraste entre deux groupes.
\end{itemize}
\end{aretenir}

\coursec{Production guidée : réemploi du lexique}

Après avoir écouté le texte, repéré les informations essentielles et appliqué la méthode d'écoute, il est temps de \cle{réemployer} activement le lexique du travail communautaire. Comprendre ne suffit pas : un bon élève de Première doit savoir \emph{produire}, à l'oral comme à l'écrit, avec les mots et les expressions de la leçon. C'est exactement ce que l'on attend de vous à l'examen : réutiliser un vocabulaire précis dans une situation de communication réelle.

\courssub{Réemploi oral : cinq phrases personnelles}

Commençons par un échauffement oral. Reprenons les cinq verbes et expressions clés de la leçon et observons d'abord des modèles tirés du script d'écoute :

\begin{exemplebox}[Modèles à imiter]
\begin{itemize}
\item \eng{Every month, we take part in communal work in our quarter.} « Chaque mois, nous participons au travail communautaire dans notre quartier. »
\item \eng{The young men clear the gutters before the rainy season.} « Les jeunes gens débouchent les caniveaux avant la saison des pluies. »
\item \eng{The women sweep the market square early in the morning.} « Les femmes balayent la place du marché tôt le matin. »
\item \eng{The council provides tools and drinking water.} « La mairie fournit les outils et l'eau potable. »
\item \eng{Everybody should join the team next Saturday.} « Tout le monde devrait rejoindre l'équipe samedi prochain. »
\end{itemize}
\end{exemplebox}

\begin{attention}[Construisez des phrases personnelles et correctes]
\begin{itemize}
\item \eng{take part in} est toujours suivi d'un nom ou d'un groupe nominal : \eng{take part in communal work}, \eng{take part in the cleaning}. Ne dites jamais \eng{take part to} (calque du français « participer à »).
\item \eng{provide} se construit ainsi : \eng{provide something} ou \eng{provide somebody with something}. Évitez le calque \eng{provide somebody something}.
\item \eng{join} ne prend pas de préposition : \eng{join the team}, \eng{join us} — jamais \eng{join to the team}.
\item \eng{clear} est régulier : \eng{clear – cleared – cleared}. \eng{sweep} est irrégulier : \eng{sweep – swept – swept} (« balayer »).
\end{itemize}
\end{attention}

\begin{experience}[À vous de parler !]
Produisez oralement \textbf{cinq phrases personnelles et vraies} (sur votre quartier, votre village ou votre école), en utilisant chacun des mots suivants : \eng{take part in}, \eng{clear}, \eng{sweep}, \eng{provide}, \eng{join}. Exemple attendu : \eng{Last Saturday, I took part in the clean-up campaign at Mokolo Market.} « Samedi dernier, j'ai participé à la campagne de nettoyage au marché Mokolo. »
\end{experience}

\courssub{Observer un modèle : l'annonce d'invitation}

Dans la vie réelle, le travail communautaire commence souvent par une \cle{annonce} (\eng{announcement}) lue à la radio locale, affichée au chef-lieu de quartier ou envoyée dans un groupe WhatsApp. Observons attentivement ce modèle avant d'en dégager la structure.

\begin{textebox}[Model announcement — communal work at Mile 4 Market]
Dear friends,

Our quarter is organising communal work this Saturday at Mile 4 Market, from 7 a.m. to 11 a.m. We shall clear the gutters and sweep the streets around the market. The health committee will provide gloves and drinking water for all the workers.

Please bring your brooms, hoes and buckets. Do not forget to wear old clothes and strong shoes. After the work, we shall share kola nuts and palm wine under the big mango tree.

Come in large numbers — together we are stronger!

Your quarter head
\end{textebox}

\textbf{Glossaire :}
\begin{itemize}
\item \eng{quarter} (n.) : quartier
\item \eng{gutters} (n. pluriel) : caniveaux
\item \eng{gloves} (n. pluriel) : gants
\item \eng{hoe} (n.) : houe, pioche
\item \eng{bucket} (n.) : seau
\item \eng{quarter head} (n.) : chef de quartier
\item \eng{kola nuts} (n. pluriel) : noix de kola
\end{itemize}

\textbf{Questions d'observation :}
\begin{enumerate}
\item Who is speaking, and to whom? (\emph{The quarter head, to the inhabitants.})
\item When and where will the work take place? (\emph{This Saturday, at Mile 4 Market, from 7 a.m. to 11 a.m.})
\item What are the two main tasks? (\emph{To clear the gutters and to sweep the streets.})
\item What must people bring? (\emph{Brooms, hoes and buckets.})
\end{enumerate}

\courssub{La structure d'une annonce : cinq informations obligatoires}

Observons l'ordre des idées dans le modèle : on salue d'abord, puis on annonce l'événement, on détaille les tâches, on précise le matériel, et on termine par un appel chaleureux. C'est un plan universel, facile à mémoriser.

\begin{popfigure}
\begin{center}
\begin{tikzpicture}[
box/.style={draw=popblue, fill=popblueL, rounded corners, thick, align=center, minimum width=3.4cm, minimum height=1cm, font=\small},
arr/.style={-{Stealth[length=3mm]}, thick, poporange}]
\node[box, align=center] (g) at (0,0){\textbf{1. Greeting}\\ \eng{Dear friends,}};
\node[box, fill=popgreenL, draw=popgreen, align=center] (w) at (4.6,0){\textbf{2. What / When / Where}\\ \eng{communal work, Saturday, Mile 4}};
\node[box, fill=poppurpleL, draw=poppurple, align=center] (t) at (9.2,0){\textbf{3. Tasks}\\ \eng{clear, sweep, plant}};
\node[box, fill=popgoldL, draw=popgold, align=center] (o) at (4.6,-2.2){\textbf{4. Tools to bring}\\ \eng{brooms, hoes, buckets}};
\node[box, fill=poppinkL, draw=poppink, align=center] (a) at (9.2,-2.2){\textbf{5. Appeal}\\ \eng{Come in large numbers!}};
\draw[arr] (g) -- (w);
\draw[arr] (w) -- (t);
\draw[arr] (t) |- (o);
\draw[arr] (o) -- (a);
\end{tikzpicture}
\end{center}
\legende{Structure d'une annonce de travail communautaire : salutation, annonce, tâches, outils, appel.}
\end{popfigure}

\begin{methode}[Rédiger une annonce en cinq étapes]
\begin{enumerate}
\item \textbf{Who} — saluez votre public : \eng{Dear friends,} / \eng{Dear neighbours,} / \eng{Dear students,}.
\item \textbf{What, when, where} — annoncez l'événement en une phrase complète : \eng{Our school is organising communal work on Friday 14th June, from 8 a.m. to noon, in the school yard.}
\item \textbf{Tasks} — énumérez deux ou trois tâches avec les verbes de la leçon : \eng{We shall sweep the classrooms, clear the grass and water the flowers.}
\item \textbf{What to bring} — précisez les outils : \eng{Please bring your cutlasses, brooms and buckets.}
\item \textbf{Appeal} — terminez par un appel à participer : \eng{Come in large numbers!} ou \eng{Everybody is welcome!}
\end{enumerate}
\end{methode}

\begin{propriete}[Les formules utiles de l'annonce]
\begin{center}
\begin{tabularx}{\linewidth}{|X|X|X|}
\hline
\rowcolor{popblueL} \textbf{Fonction} & \textbf{English} & \textbf{Français} \\
\hline
Salutation & \eng{Dear friends, / Dear neighbours,} & Chers amis, / Chers voisins, \\
\hline
Annoncer & \eng{Our quarter is organising\ldots} & Notre quartier organise\ldots \\
\hline
Dater & \eng{this Saturday, on 14th June} & ce samedi, le 14 juin \\
\hline
Situer & \eng{at Mile 4 Market, in the school yard} & au marché Mile 4, dans la cour de l'école \\
\hline
Donner l'heure & \eng{from 7 a.m. to 11 a.m.} & de 7 h à 11 h \\
\hline
Demander & \eng{Please bring your brooms.} & Apportez vos balais, s'il vous plaît. \\
\hline
Appeler & \eng{Come in large numbers!} & Venez nombreux ! \\
\hline
Conclure & \eng{Together we are stronger!} & Ensemble, nous sommes plus forts ! \\
\hline
\end{tabularx}
\end{center}
\end{propriete}

\begin{attention}[Pièges de traduction pour les francophones]
\begin{itemize}
\item Ne traduisez jamais « venez nombreux » par \eng{come numerous} : c'est un calque incorrect. Dites \eng{Come in large numbers!} ou \eng{Everybody is welcome!}
\item « de 7 h à 11 h » se dit \eng{from 7 a.m. to 11 a.m.} — n'oubliez ni \eng{from} ni \eng{to}.
\item \eng{Please bring your brooms} signifie « Apportez vos balais, s'il vous plaît ». \eng{bring} (« apporter », vers le locuteur) ne se confond pas avec \eng{take} (« emporter », loin du locuteur).
\item En anglais, on écrit \eng{on Saturday} (et non \eng{in Saturday}) et \eng{at 7 o'clock} (et non \eng{in 7 o'clock}).
\end{itemize}
\end{attention}

\begin{exoresolu}[Rédiger une annonce complète]
Énoncé : Your school, Government High School Buea, is organising communal work on Saturday 21st September, from 7 a.m. to 10 a.m. Tasks: sweep the classrooms, clear the grass around the buildings and plant flowers. Students must bring cutlasses, brooms and buckets. Write a short announcement (about 60 words) inviting all students to take part.
\tcblower
Solution détaillée : on suit le plan en cinq étapes.

\eng{Dear students,}

\eng{Our school is organising communal work on Saturday 21st September, from 7 a.m. to 10 a.m., in the school compound. We shall sweep the classrooms, clear the grass around the buildings and plant flowers. Please bring your cutlasses, brooms and buckets. Come in large numbers — a clean school is a healthy school!}

\eng{The Principal}

Vérification : salutation (1), what/when/where (2), trois tâches avec \eng{sweep}, \eng{clear}, \eng{plant} (3), outils avec \eng{bring} (4), appel final (5). Toutes les informations obligatoires sont présentes, en environ 60 mots.
\end{exoresolu}

\courssub{Expression orale en binômes et correction collective}

\begin{experience}[Pair work : read, take notes, ask questions]
Travaillez en binômes.
\begin{enumerate}
\item L'élève A rédige puis \textbf{lit à voix haute} son annonce (travail communautaire dans son quartier, son village ou son école).
\item L'élève B \textbf{prend des notes} pendant l'écoute : date, lieu, heure, tâches, outils — exactement comme dans la méthode d'écoute de la section précédente.
\item L'élève B pose ensuite \textbf{deux questions} pour vérifier ses notes, par exemple : \eng{What time does the work start?} « À quelle heure le travail commence-t-il ? » ; \eng{What must we bring?} « Que devons-nous apporter ? »
\item Échangez les rôles.
\end{enumerate}
\end{experience}

\begin{exercice}[Correction collective au tableau]
Deux binômes volontaires écrivent leur annonce au tableau. Avec toute la classe, corrigez-les en vérifiant les cinq points suivants :
\begin{enumerate}
\item Les cinq informations obligatoires sont-elles présentes (who, what, when, where, what to bring) ?
\item Les prépositions sont-elles correctes (\eng{on Saturday}, \eng{at 7 a.m.}, \eng{from\ldots to\ldots}) ?
\item Les verbes de la leçon sont-ils bien employés (\eng{take part in}, \eng{clear}, \eng{sweep}, \eng{provide}, \eng{join}) ?
\item L'appel final est-il idiomatique (\eng{Come in large numbers!} et non \eng{come numerous}) ?
\item L'orthographe et la ponctuation sont-elles soignées ?
\end{enumerate}
\end{exercice}

\begin{corrige}[Grille de correction type]
Une annonce réussie ressemble à ceci : \eng{Dear neighbours, our village is organising communal work this Saturday at the community hall, from 6.30 a.m. to 10 a.m. We shall clear the gutters, sweep the main road and cut the grass. The council will provide drinking water. Please bring your hoes, cutlasses and wheelbarrows. Everybody is welcome — together we are stronger!} Erreurs typiques à corriger au tableau : \eng{in Saturday} $\rightarrow$ \eng{on Saturday} ; \eng{take part to} $\rightarrow$ \eng{take part in} ; \eng{come numerous} $\rightarrow$ \eng{come in large numbers} ; \eng{join to us} $\rightarrow$ \eng{join us}.
\end{corrige}

\begin{savaistu}[Le travail communautaire au Cameroun et dans le Commonwealth]
Au Cameroun anglophone (régions du Nord-Ouest et du Sud-Ouest), le \eng{communal work} (ou \eng{community labour}) est une tradition forte. Chaque \eng{quarter} (quartier) ou village organise des \eng{clean-up campaigns} le \eng{Country Sunday} (souvent le dernier samedi du mois) ou sur convocation du \eng{quarter head} (chef de quartier) ou du \eng{Traditional Council}. Les tâches typiques sont : \eng{clear the gutters}, \eng{sweep the streets}, \eng{cut the grass}, \eng{repair the road}. La mairie (\eng{the council}) fournit parfois l'eau (\eng{drinking water}) et les outils (\eng{tools : cutlasses, hoes, wheelbarrows}). Au Kenya, on parle de \eng{Harambee} (« tous ensemble » en swahili) ; au Royaume-Uni, de \eng{volunteering} ou \eng{community service}. L'esprit est le même : \eng{Together we are stronger!}
\end{savaistu}

\begin{aretenir}[L'essentiel de la production guidée]
Pour réemployer le lexique du travail communautaire : (1) construisez des phrases personnelles avec \eng{take part in}, \eng{clear}, \eng{sweep}, \eng{provide}, \eng{join} en respectant leurs constructions ; (2) une annonce suit toujours le plan \textbf{Greeting $\rightarrow$ What/When/Where $\rightarrow$ Tasks $\rightarrow$ Tools $\rightarrow$ Appeal} ; (3) les cinq informations obligatoires sont \textbf{who, what, when, where, what to bring} ; (4) méfiez-vous des calques du français : \eng{come in large numbers}, \eng{on Saturday}, \eng{from 7 a.m. to 11 a.m.}.
\end{aretenir}

\newpage

\coursec{Activité d'intégration}

\coursec{Lesson 4 — Listening: Texts about taking part in communal work}

\courssub{Objectifs de l'activité}

À l'issue de cette activité d'intégration, tu seras capable de :

\begin{itemize}
\item \cle{réemployer} le lexique du travail communautaire (\eng{communal work, clean-up, tools, volunteer, take part in}) dans une production orale authentique ;
\item \cle{produire} une annonce radio claire et convaincante de 8 à 10 lignes, avec les informations essentielles : lieu, date, heures, tâches, matériel ;
\item \cle{appliquer} les stratégies d'écoute étudiées dans la leçon : repérer le thème, noter les mots-clés, les chiffres et les noms propres pendant l'écoute des autres groupes ;
\item \cle{poser} des questions de compréhension en anglais correct (\eng{Wh-} questions) ;
\item \cle{évaluer} une production orale selon des critères simples : clarté, exactitude des informations, pouvoir de conviction.
\end{itemize}

\courssub{Situation}

\begin{textebox}[Contexte de la tâche]
You are members of the youth club of your neighbourhood in Buea. The quarter is dirty after the rainy season: the gutters are blocked, plastic bags are everywhere and the community market needs cleaning. The quarter head has asked your youth club to organise a communal work day next Saturday. Your club has been offered five minutes on the local radio station, \emph{Mount Cameroon Radio}, to invite the population. Prepare your radio announcement!
\end{textebox}

Traduction du contexte : vous êtes membres du club des jeunes de votre quartier à Buéa. Le quartier est sale après la saison des pluies : les caniveaux sont bouchés, les sachets plastiques sont partout et le marché communautaire a besoin d'être nettoyé. Le chef de quartier a demandé à votre club d'organiser une journée de travail communautaire samedi prochain. Votre club dispose de cinq minutes à la radio locale, \emph{Mount Cameroon Radio}, pour inviter la population. Préparez votre annonce !

\begin{definition}[Rappel : qu'est-ce qu'une annonce radio efficace ?]
Une \cle{annonce radio} (\eng{radio announcement}) est un message oral court qui doit : attirer l'attention dès la première phrase, donner les informations essentielles (\eng{who, what, where, when, why}), être répétée dans les points clés (l'auditeur ne peut pas revenir en arrière), et se terminer par un \cle{appel à l'action} (\eng{call to action}).
\end{definition}

\courssub{Consignes}

Travail par groupes de 4 élèves. Durée totale : 50 minutes.

\begin{enumerate}
\item \textbf{Brainstorming (5 min).} Dans votre groupe, listez en anglais : le nom de votre quartier ou village, la date et les heures du travail communautaire, \cle{trois tâches précises} (ex. \eng{clear the gutters, sweep the market square, collect plastic bags}), et les \cle{outils à apporter} (ex. \eng{cutlasses, brooms, rakes, wheelbarrows, gloves}).
\item \textbf{Rédaction (15 min).} Rédigez ensemble une annonce de \cle{8 à 10 lignes}. Elle doit contenir : une accroche (\eng{Hello listeners! / Attention please!}), le nom du lieu, la date et les heures, les trois tâches, le matériel à apporter, et un appel final convaincant. Utilisez le \cle{futur} (\eng{will}) et l'\cle{impératif} (\eng{Come and join us! Bring your tools!}).
\item \textbf{Répétition (5 min).} Distribuez les rôles : un présentateur, un ou deux voix qui donnent les détails, une voix pour l'appel final. Répétez à voix haute en veillant à la prononciation et à l'intonation.
\item \textbf{Présentation et écoute active (15 min).} Chaque groupe présente son annonce devant la classe (ou joue son enregistrement). Pendant l'écoute, les autres groupes remplissent la \cle{fiche d'écoute} ci-dessous, puis posent \cle{deux questions de compréhension} au groupe qui vient de présenter.
\item \textbf{Vote (5 min).} La classe vote pour l'annonce la plus claire et la plus convaincante, en justifiant son choix en anglais : \eng{I vote for group 3 because their announcement was clear and complete.}
\end{enumerate}

\begin{exemplebox}[Fiche d'écoute à compléter pendant les présentations]
\begin{center}
\begin{tabularx}{\linewidth}{|X|X|}\hline
\rowcolor{popblueL} \textbf{Question} & \textbf{Your notes} \\ \hline
\eng{Where is the communal work? (place)} & \\ \hline
\eng{When is it? (day, date, time)} & \\ \hline
\eng{What are the three tasks?} & 1. \quad 2. \quad 3. \\ \hline
\eng{What tools must people bring?} & \\ \hline
\eng{Who organises the event?} & \\ \hline
\end{tabularx}
\end{center}
\end{exemplebox}

\courssub{Ressources à mobiliser}

\begin{aretenir}[Lexique utile]
\begin{itemize}
\item \eng{communal work / community work} : travail communautaire ; \eng{a clean-up campaign} : une campagne de nettoyage ;
\item \eng{to take part in / to join / to volunteer} : participer à / se joindre à / se porter volontaire ;
\item tâches : \eng{to clear the gutters} (déboucher les caniveaux), \eng{to sweep the streets} (balayer les rues), \eng{to collect rubbish} (ramasser les ordures), \eng{to cut the grass} (couper l'herbe), \eng{to plant flowers} (planter des fleurs) ;
\item outils : \eng{a cutlass} (une machette), \eng{a broom} (un balai), \eng{a rake} (un râteau), \eng{a wheelbarrow} (une brouette), \eng{a hoe} (une houe), \eng{gloves} (des gants), \eng{a bucket} (un seau).
\end{itemize}
\end{aretenir}

\begin{propriete}[Structures grammaticales à réemployer]
\begin{itemize}
\item Le \cle{futur avec will} pour annoncer l'événement : \eng{The communal work will take place on Saturday.} « Le travail communautaire aura lieu samedi. »
\item L'\cle{impératif} pour l'appel à l'action : \eng{Come in large numbers! Bring your brooms! Don't miss it!} « Venez nombreux ! Apportez vos balais ! Ne le manquez pas ! »
\item Les \cle{questions en Wh-} pour la compréhension : \eng{When will the work start? What tools must we bring? Who is the organiser?}
\item Les expressions de temps : \eng{from 7 a.m. to 12 noon} (de 7 h à midi), \eng{next Saturday} (samedi prochain), \eng{in front of the market} (devant le marché).
\end{itemize}
\end{propriete}

\begin{attention}[Erreurs fréquentes à éviter]
\begin{itemize}
\item Ne dis pas \eng{*a communitary work} : dis \eng{communal work} ou \eng{community work}.
\item \eng{To assist} est un faux ami : il signifie « aider », pas « assister à ». Pour « assister à », utilise \eng{to attend} ou \eng{to take part in}.
\item Heure : dis \eng{at 7 o'clock} ou \eng{from 7 a.m. to noon}, jamais \eng{*at 7 hours}.
\item N'oublie pas la majuscule aux jours et mois : \eng{Saturday, June}.
\end{itemize}
\end{attention}

\courssub{Proposition de production et corrigé commenté}

\begin{exoresolu}[Modèle d'annonce radio]
Énoncé : rédigez l'annonce radio de votre groupe (8 à 10 lignes) pour inviter la population au travail communautaire de samedi prochain.
\tcblower
\textbf{Production modèle :}

“Hello listeners! This is Mount Cameroon Radio, the voice of your neighbourhood. The youth club of Bonduma quarter invites you to a big communal work day next Saturday, 15th June. The work will take place in front of the community market, from 7 o'clock in the morning to 12 noon. Together, we will clear the blocked gutters, sweep the market square and collect all the plastic bags in the streets. Please bring your cutlasses, brooms, rakes and wheelbarrows. Don't forget your gloves! Come in large numbers, because a clean quarter is a healthy quarter. Join us this Saturday — your quarter needs you!”

« Bonjour chers auditeurs ! Ici Mount Cameroon Radio, la voix de votre quartier. Le club des jeunes du quartier Bonduma vous invite à une grande journée de travail communautaire samedi prochain, le 15 juin. Le travail aura lieu devant le marché communautaire, de 7 heures du matin à midi. Ensemble, nous déboucherons les caniveaux obstrués, balayerons la place du marché et ramasserons tous les sachets plastiques dans les rues. Apportez vos machettes, balais, râteaux et brouettes, s'il vous plaît. N'oubliez pas vos gants ! Venez nombreux, car un quartier propre est un quartier sain. Rejoignez-nous ce samedi — votre quartier a besoin de vous ! »

\textbf{Commentaire des choix de langue :}
\begin{itemize}
\item \cle{Accroche} : \eng{Hello listeners!} capte immédiatement l'attention, technique typique des annonces radio.
\item \cle{Informations essentielles} dès la deuxième phrase : qui (\eng{the youth club of Bonduma}), quoi (\eng{a big communal work day}), quand (\eng{next Saturday, 15th June}).
\item \cle{Futur en will} pour les tâches : \eng{we will clear, sweep and collect} — trois verbes d'action précis, conformes à la consigne.
\item \cle{Impératifs} pour l'appel à l'action : \eng{bring, don't forget, come, join} — ton direct et motivant.
\item \cle{Slogan final} : \eng{a clean quarter is a healthy quarter} — phrase courte et mémorable, efficace à la radio.
\item Répétition du jour à la fin (\eng{this Saturday}) : l'auditeur retient l'information clé.
\end{itemize}
\end{exoresolu}

\courssub{Bilan de l'activité}

\begin{aretenir}[Ce que cette activité t'a appris]
\begin{itemize}
\item \textbf{Compréhension orale} : tu as écouté des annonces authentiques produites par tes camarades et tu as relevé les informations essentielles (lieu, date, tâches, outils) grâce aux stratégies de la leçon : repérage des mots-clés, des chiffres et des noms propres.
\item \textbf{Expression orale} : tu as produit un message communicatif réel, avec une intention claire (inviter, convaincre), en soignant prononciation et intonation.
\item \textbf{Lexique} : tu as réemployé le vocabulaire du travail communautaire dans un contexte naturel et camerounais.
\item \textbf{Grammaire} : tu as mobilisé le futur en \eng{will}, l'impératif et les questions en \eng{Wh-}.
\item \textbf{Compétence sociale} : tu as travaillé en équipe, évalué tes pairs et justifié ton choix en anglais — exactement l'esprit du travail communautaire !
\end{itemize}
\end{aretenir}

\begin{experience}[Pour aller plus loin]
À la maison, écoute une vraie annonce à la radio (CRTVC, une radio communautaire ou une station anglophone) et remplis la même fiche d'écoute : \eng{where, when, tasks, tools}. Au cours suivant, présente tes réponses en trois phrases : \eng{I listened to an announcement about... It will take place... People must bring...}
\end{experience}

\newpage

\coursec{Méthodes et exercices résolus}

\courssub{Fiche méthode 1 — Comprendre un enregistrement : repérer le thème et les mots-clés}

\begin{textebox}[Listening script — « The Bafoussam clean-up day »]
Every last Saturday of the month, the people of Bafoussam take part in communal work. At seven o'clock in the morning, about two hundred volunteers meet in front of the town hall. The mayor, Mr Kamga, gives them tools: brooms, wheelbarrows and gloves. The young people clean the market square, while the women sweep the streets near the health centre. The work lasts three hours. After that, the council offers food and drinks to everybody. Mr Kamga says that communal work keeps the town clean and brings people together.
\end{textebox}

\textbf{Glossaire :} \eng{communal work} (nom, travail communautaire) ; \eng{a volunteer} (nom, un bénévole) ; \eng{the town hall} (nom, la mairie) ; \eng{a broom} (nom, un balai) ; \eng{a wheelbarrow} (nom, une brouette) ; \eng{to sweep} (verbe, balayer — prétérit \eng{swept}) ; \eng{the council} (nom, le conseil municipal) ; \eng{to bring people together} (expression, rassembler les gens).

\begin{methode}[Réussir une écoute sur le travail communautaire]
\begin{enumerate}
\item \textbf{Avant l'écoute} : lis d'abord toutes les questions. Elles annoncent le thème et te disent quoi écouter. Souligne les mots interrogatifs (\eng{who, where, when, how many, why}).
\item \textbf{Première écoute} : ne note rien. Identifie seulement le \cle{thème général} (ici : \eng{communal work}, \eng{community clean-up}, \eng{village development}) et les \cle{personnes} qui parlent.
\item \textbf{Deuxième écoute} : note les \cle{mots-clés} : noms propres (\eng{the mayor, Mr Kamga}), lieux (\eng{the market square, the health centre}), chiffres (\eng{two hundred volunteers, three hours}), jours et heures (\eng{every last Saturday, at seven o'clock}).
\item \textbf{Après l'écoute} : réponds en phrases courtes et simples. Reprends les mots de la question : \eng{When does the work start?} $\rightarrow$ \eng{The work starts at seven o'clock.}
\item \textbf{Vérification} : relis chaque réponse : sujet + verbe conjugué, majuscule aux noms propres, préposition correcte (\eng{at} + heure, \eng{on} + jour).
\end{enumerate}
\end{methode}

\begin{exoresolu}[Listening comprehension — « The Bafoussam clean-up day »]
\textbf{Script lu deux fois par l'enseignant :}

\eng{Every last Saturday of the month, the people of Bafoussam take part in communal work. At seven o'clock in the morning, about two hundred volunteers meet in front of the town hall. The mayor, Mr Kamga, gives them tools: brooms, wheelbarrows and gloves. The young people clean the market square, while the women sweep the streets near the health centre. The work lasts three hours. After that, the council offers food and drinks to everybody. Mr Kamga says that communal work keeps the town clean and brings people together.}

\textbf{Questions :} 1. How often do the people of Bafoussam do communal work? 2. Where do the volunteers meet? 3. How many volunteers are there? 4. What do the young people clean? 5. Give two benefits of communal work according to the mayor.
\tcblower
\textbf{Méthode appliquée :} avant l'écoute, on repère les mots interrogatifs : \eng{how often} (fréquence), \eng{where} (lieu), \eng{how many} (nombre), \eng{what} (objet), \eng{benefits} (avantages). Pendant l'écoute, on note : \eng{last Saturday / town hall / two hundred / market square / clean + brings together}.

\textbf{Réponses rédigées :}
\begin{enumerate}
\item \eng{They do communal work once a month, every last Saturday.} (\eng{every last Saturday of the month} = une fois par mois.)
\item \eng{The volunteers meet in front of the town hall.} (Préposition \eng{in front of} = devant ; ne pas confondre avec \eng{in} = à l'intérieur.)
\item \eng{There are about two hundred volunteers.} (\eng{about} = environ ; attention, \eng{hundred} reste invariable après un chiffre : \eng{two hundred}, jamais \emph{two hundreds}.)
\item \eng{The young people clean the market square.}
\item \eng{According to the mayor, communal work keeps the town clean and brings people together.} (On reformule avec \eng{according to} = selon.)
\end{enumerate}
\textbf{Conclusion :} chaque réponse reprend les mots de la question, avec un verbe conjugué au présent simple (habitude : \eng{take part, meet, clean, lasts}).
\end{exoresolu}

\courssub{Fiche méthode 2 — Répondre aux questions vrai/faux et aux QCM}

\begin{methode}[Vrai/Faux et QCM : justifier par le texte]
\begin{enumerate}
\item Lis chaque affirmation \textbf{entièrement} avant l'écoute : une seule information fausse rend toute la phrase fausse.
\item Repère le \cle{piège} possible : un chiffre changé (\eng{fifty} au lieu de \eng{two hundred}), un lieu changé, une négation ajoutée.
\item Pendant l'écoute, coche mentalement chaque détail entendu ; méfie-toi des mots proches : \eng{council} (conseil municipal) et \eng{counsel} (conseil, avis) se prononcent pareil.
\item Pour le QCM, élimine d'abord les options impossibles, puis choisis celle qui reprend les mots exacts du script.
\item \textbf{Justifie toujours} : cite la phrase du script qui prouve ta réponse (\eng{The text says: “The work lasts three hours.”}).
\end{enumerate}
\end{methode}

\begin{exoresolu}[True or False + QCM — même script « The Bafoussam clean-up day »]
\textbf{A. True or False. Justify with a sentence from the text.}
\begin{enumerate}
\item \eng{Communal work takes place every week.}
\item \eng{The mayor gives tools to the volunteers.}
\item \eng{The women clean the market square.}
\end{enumerate}
\textbf{B. Choose the correct answer.}

\eng{The work lasts: a) two hours \quad b) three hours \quad c) the whole day}
\tcblower
\textbf{A. True or False :}
\begin{enumerate}
\item \textbf{False.} \eng{The text says: “Every last Saturday of the month, the people of Bafoussam take part in communal work.”} Le travail a lieu une fois par mois, pas chaque semaine. Piège classique : \eng{every week} $\neq$ \eng{every last Saturday of the month}.
\item \textbf{True.} \eng{The text says: “The mayor, Mr Kamga, gives them tools: brooms, wheelbarrows and gloves.”} Information exactement conforme au script.
\item \textbf{False.} \eng{The text says: “The young people clean the market square, while the women sweep the streets near the health centre.”} Les rôles ont été inversés : ce sont les jeunes qui nettoient la place du marché.
\end{enumerate}
\textbf{B. QCM :} la bonne réponse est \textbf{b) three hours}. Justification : \eng{The text says: “The work lasts three hours.”} L'option a) correspond à un chiffre proche (piège fréquent à l'écoute entre \eng{two} et \eng{three}) ; l'option c) exagère la durée.

\textbf{Conclusion :} en vrai/faux comme en QCM, la justification par une citation exacte du script est obligatoire pour obtenir tous les points.
\end{exoresolu}

\courssub{Fiche méthode 3 — Relever et exploiter le lexique du travail communautaire}

\begin{methode}[Construire et réemployer le vocabulaire thématique]
\begin{enumerate}
\item Classe le vocabulaire en \cle{familles} : les personnes (\eng{volunteers, the mayor, the community leader}), les outils (\eng{a broom, a wheelbarrow, a rake, gloves, a machete}), les actions (\eng{to sweep, to clean, to plant, to dig, to collect rubbish}), les lieux (\eng{the market square, the health centre, the village square}).
\item Apprends chaque mot avec son \cle{article} et sa \cle{préposition} : \eng{to take part in} (participer à), \eng{to work together}, \eng{to keep the town clean}.
\item Note les \cle{expressions utiles} : \eng{communal work} (travail communautaire), \eng{clean-up campaign} (campagne de nettoyage), \eng{to bring people together} (rassembler les gens).
\item Réemploie chaque mot dans une \textbf{phrase personnelle} au présent simple (habitude) ou au prétérit (événement passé).
\item Vérifie l'orthographe des mots pièges : \eng{volunteer} (un seul \emph{l}), \eng{government}, \eng{environment}.
\end{enumerate}
\end{methode}

\begin{definition}[Lexique essentiel du travail communautaire]
\begin{center}
\begin{tabularx}{\linewidth}{|X|X|X|}
\hline
\rowcolor{popblueL} English & Français & Remarque \\
\hline
communal work & le travail communautaire & nom indénombrable, sans article \\
\hline
a volunteer & un bénévole & un seul \emph{l} ; prononcé « vo-lon-tire » \\
\hline
a broom / a rake & un balai / un râteau & outils de nettoyage \\
\hline
a wheelbarrow & une brouette & prononcé « ouil-bé-ro » \\
\hline
rubbish & les ordures & indénombrable : jamais \emph{rubbishes} \\
\hline
to sweep (swept, swept) & balayer & verbe irrégulier \\
\hline
to take part in & participer à & préposition \eng{in} obligatoire \\
\hline
the community leader & le chef de la communauté & aussi \eng{the village chief} \\
\hline
a clean-up campaign & une campagne de nettoyage & \eng{clean-up} avec trait d'union \\
\hline
to bring people together & rassembler les gens & verbe irrégulier : bring — brought — brought \\
\hline
\end{tabularx}
\end{center}
\end{definition}

\begin{exoresolu}[Vocabulary — matching et réemploi]
\textbf{A. Match each word with its French equivalent.}

\begin{center}
\begin{tabularx}{\linewidth}{|X|X|}
\hline
\rowcolor{popblueL} English & Français \\
\hline
1. a wheelbarrow & a) un balai \\
\hline
2. a broom & b) les ordures \\
\hline
3. rubbish & c) une brouette \\
\hline
4. a volunteer & d) un bénévole \\
\hline
\end{tabularx}
\end{center}

\textbf{B. Complete the sentences with:} \eng{takes part in — tools — together — sweep}

1. \eng{The villagers \ldots\ the streets every Saturday.}
2. \eng{Everybody \ldots\ communal work in my village.}
3. \eng{The council provides \ldots\ such as rakes and gloves.}
4. \eng{Communal work brings people \ldots .}
\tcblower
\textbf{A. Correspondances :} 1 $\rightarrow$ c) (\eng{a wheelbarrow} = une brouette) ; 2 $\rightarrow$ a) (\eng{a broom} = un balai) ; 3 $\rightarrow$ b) (\eng{rubbish} = les ordures, nom indénombrable : jamais \emph{rubbishes}) ; 4 $\rightarrow$ d) (\eng{a volunteer} = un bénévole).

\textbf{B. Phrases complétées :}
\begin{enumerate}
\item \eng{The villagers \textbf{sweep} the streets every Saturday.} (Présent simple, sujet pluriel : verbe sans \emph{s}.)
\item \eng{Everybody \textbf{takes part in} communal work in my village.} (\eng{everybody} est singulier : \eng{takes} avec \emph{s} ; expression figée \eng{to take part in} + nom.)
\item \eng{The council provides \textbf{tools} such as rakes and gloves.} (\eng{such as} = comme, par exemple.)
\item \eng{Communal work brings people \textbf{together}.} (Expression : \eng{to bring people together} = rassembler.)
\end{enumerate}
\textbf{Conclusion :} le lexique s'apprend en contexte ; chaque mot réemployé dans une phrase correcte vaut des points à l'examen.
\end{exoresolu}

\courssub{Fiche méthode 4 — Relever chiffres, noms propres et détails}

\begin{methode}[Noter les informations chiffrées et les noms propres]
\begin{enumerate}
\item Prépare avant l'écoute une petite grille mentale : \cle{qui} (names), \cle{où} (places), \cle{quand} (days, times), \cle{combien} (numbers).
\item Écris les chiffres en \textbf{lettres} dans ta réponse si le nombre est petit (\eng{three hours, fifty volunteers}) : c'est la norme en rédaction anglaise.
\item Attention aux confusions sonores : \eng{thirteen/thirty}, \eng{fifteen/fifty} — la terminaison \emph{-teen} est accentuée, \emph{-ty} ne l'est pas.
\item Les noms propres prennent toujours une \cle{majuscule} : \eng{Mr Kamga, Bafoussam, Saturday}. Pas d'article devant les villes.
\item Restitue le détail avec la bonne préposition : \eng{at seven o'clock}, \eng{on Saturday}, \eng{in the morning}.
\end{enumerate}
\end{methode}

\begin{exoresolu}[Short answers — détails précis]
\textbf{Script (extrait lu deux fois) :}

\eng{In Bamenda, communal work starts at six thirty in the morning and finishes at nine. Last month, one hundred and fifty volunteers cleaned the commercial avenue. Mrs Abena, the community leader, thanked everybody and announced that the next clean-up day will take place on the twenty-fifth of June.}

\textbf{Questions :} 1. What time does communal work start in Bamenda? 2. How long does it last? 3. How many volunteers cleaned the avenue last month? 4. Who is Mrs Abena? 5. When will the next clean-up day take place?
\tcblower
\textbf{Repérage pendant l'écoute :} heure = \eng{six thirty} ; fin = \eng{nine} ; nombre = \eng{one hundred and fifty} ; personne = \eng{Mrs Abena, community leader} ; date = \eng{the twenty-fifth of June}.

\textbf{Réponses rédigées :}
\begin{enumerate}
\item \eng{It starts at six thirty in the morning.} (Préposition \eng{at} devant l'heure.)
\item \eng{It lasts two and a half hours.} (De 6 h 30 à 9 h = 2 h 30 ; on dit \eng{two and a half hours} ou \eng{two hours and a half}, mais jamais \emph{two and half hours}.)
\item \eng{One hundred and fifty volunteers cleaned the avenue last month.} (Prétérit : \eng{cleaned}, car \eng{last month} = événement passé ; \eng{hundred} invariable.)
\item \eng{Mrs Abena is the community leader.} (Majuscules à \eng{Mrs Abena} ; pas d'article devant le nom propre.)
\item \eng{The next clean-up day will take place on the twenty-fifth of June.} (Futur avec \eng{will} car c'est une annonce ; \eng{on} devant la date.)
\end{enumerate}
\textbf{Conclusion :} chaque détail est restitué avec la préposition et le temps verbal corrects : présent pour l'habitude, prétérit pour le passé, \eng{will} pour l'annonce future.
\end{exoresolu}

\courssub{Fiche méthode 5 — Production guidée : réemployer le lexique à l'oral et à l'écrit}

\begin{methode}[Produire cinq phrases sur le travail communautaire]
\begin{enumerate}
\item Choisis une \cle{structure simple} : Sujet + verbe + complément + (préposition + lieu/temps).
\item Varie les temps selon le sens : présent simple pour les habitudes (\eng{We take part in communal work every month.}), prétérit pour un événement passé (\eng{Last Saturday, we cleaned the school yard.}), \eng{will} ou \eng{going to} pour un projet.
\item Intègre au moins \textbf{trois mots du lexique} de la leçon dans chaque production : \eng{volunteers, tools, to sweep, rubbish, community leader}.
\item Ajoute une phrase d'opinion avec \eng{because} : \eng{I like communal work because it brings people together.}
\item Relis-toi : accord sujet-verbe (\eng{he takes}), majuscules, prépositions (\eng{in the morning}, \eng{on Saturday}).
\end{enumerate}
\end{methode}

\begin{exoresolu}[Guided production — « Communal work in my village »]
\textbf{Énoncé :} Write five sentences about communal work in your village or town. Use at least five words from the lesson (\eng{volunteers, tools, to sweep, rubbish, to take part in, community leader, together}) and three different tenses.
\tcblower
\textbf{Plan :} phrase 1 = habitude (présent) ; phrase 2 = outils (présent) ; phrase 3 = événement passé (prétérit) ; phrase 4 = projet (futur) ; phrase 5 = opinion avec \eng{because}.

\textbf{Production modèle :}

\eng{In my village, everybody takes part in communal work on the last Saturday of the month. The community leader gives us tools such as brooms and wheelbarrows. Last month, fifty volunteers swept the streets and collected the rubbish near the market. Next Saturday, we are going to clean the area around the health centre. I enjoy communal work because it keeps our village clean and brings people together.}

\textbf{Justifications grammaticales :}
\begin{itemize}
\item \eng{takes part in} : présent simple, \eng{everybody} est singulier, d'où le \emph{s}.
\item \eng{gives} : présent simple, troisième personne du singulier.
\item \eng{swept} : prétérit du verbe irrégulier \eng{to sweep} (sweep — swept — swept) ; justifié par \eng{last month}.
\item \eng{are going to clean} : futur proche pour un projet déjà décidé (\eng{next Saturday}).
\item \eng{because it keeps\ldots and brings\ldots} : deux verbes coordonnés, chacun avec son \emph{s} de la troisième personne.
\end{itemize}
\textbf{Conclusion :} cinq phrases, cinq mots du lexique, trois temps différents : la consigne est entièrement respectée.
\end{exoresolu}

\begin{experience}[Speaking — Interview : « Our clean-up day »]
Par groupes de deux, jouez une interview entre un journaliste de la radio locale et un \eng{community leader}. Le journaliste pose au moins quatre questions avec \eng{when, where, how many, why} ; le chef répond en phrases complètes avec le lexique de la leçon. Exemple de départ : \eng{When does communal work take place in your village?} — \eng{It takes place every last Saturday of the month, at seven o'clock in the morning.} Puis inversez les rôles.
\end{experience}

\begin{savaistu}[Le travail communautaire au Cameroun et dans le monde anglophone]
Au Cameroun, le travail communautaire est une tradition vivante : les habitants d'un quartier ou d'un village se réunissent pour balayer les rues, déboucher les caniveaux ou construire une école. Dans les pays anglophones, on parle de \eng{community service} ou de \eng{volunteering} : aux États-Unis et au Kenya, les \eng{clean-up days} rassemblent aussi des centaines de \eng{volunteers}. Le mot anglais \eng{volunteer} vient du latin \emph{voluntas}, « la volonté, le bon vouloir » : un bénévole est celui qui agit de son plein gré.
\end{savaistu}

\begin{attention}[Les 5 erreurs qui coûtent des points au Bac]
\begin{enumerate}
\item \textbf{Faux ami \eng{to assist}} : \eng{to assist} signifie « aider », pas « assister à ». Pour « participer à », dis \eng{to take part in} ou \eng{to attend}. \eng{We take part in communal work.} (jamais \emph{we assist communal work}).
\item \textbf{\eng{hundred} avec un \emph{s}} : après un chiffre, \eng{hundred} est invariable : \eng{two hundred volunteers}. Le \emph{s} n'apparaît que dans \eng{hundreds of people} (des centaines de personnes).
\item \textbf{Calque du français « prendre part »} : ne dis pas \emph{take part of} ; la préposition correcte est \eng{in} : \eng{to take part \textbf{in} communal work}. De même, « participer » ne se traduit jamais par \emph{to participate to}.
\item \textbf{Prépositions de temps} : \eng{at} + heure (\eng{at seven o'clock}), \eng{on} + jour et date (\eng{on Saturday, on the twenty-fifth of June}), \eng{in} + mois et moment de la journée (\eng{in June, in the morning}). Ne les mélange pas.
\item \textbf{Oubli du \emph{s} de la troisième personne au présent simple} : \eng{Everybody take\textbf{s} part in communal work.} \eng{The work last\textbf{s} three hours.} C'est la faute la plus fréquente des candidats francophones.
\end{enumerate}
\end{attention}

\newpage

\coursec{Exercices}

\courssub{Je vérifie mes connaissances}

\begin{exercice}[QCM de compréhension]
Écoutez (ou lisez) le dialogue suivant entre un \eng{quarter head} (chef de quartier) et un journaliste, puis choisissez la bonne réponse (a, b ou c).
\begin{textebox}[A radio interview in Bamenda]
Journalist: Good morning, Mr Ndi. You are the quarter head of Ntarikon. Can you tell us about the communal work you organised last Saturday?
Quarter head: Good morning. Yes, we cleaned the main market and cleared the gutters around it.
Journalist: How many people took part?
Quarter head: About eighty volunteers, mostly young people and women from the neighbourhood.
Journalist: When did you start?
Quarter head: We started at six o'clock in the morning and finished around noon.
Journalist: Why is this work so important to you?
Quarter head: Because a clean market means fewer diseases, and it brings the community together.
\end{textebox}
\begin{enumerate}
\item Who is Mr Ndi? \quad a) a journalist \quad b) the quarter head \quad c) a market seller
\item What did the volunteers clean? \quad a) the school \quad b) the bridge \quad c) the market and the gutters
\item How many volunteers took part? \quad a) eight \quad b) eighteen \quad c) about eighty
\item When did the work start? \quad a) at 6 a.m. \quad b) at noon \quad c) at 6 p.m.
\item Why is communal work important to Mr Ndi? \quad a) it is paid \quad b) it reduces diseases and unites people \quad c) it pleases the radio
\end{enumerate}
\end{exercice}

\begin{exercice}[Vrai ou faux justifié]
Écoutez (ou lisez) le script suivant, puis dites si les affirmations sont \eng{true} ou \eng{false}. Justifiez chaque réponse par une courte citation du texte.
\begin{textebox}[Community clean-up in Buea]
Last month, the people of Molyko organised a communal work session to clean the streets after the rainy season. More than one hundred and twenty inhabitants came with their own tools: hoes, cutlasses, wheelbarrows and brooms. The work lasted four hours, from seven to eleven in the morning. The council provided gloves and drinking water for everyone. At the end, the quarter head thanked the volunteers and announced that the next session would take place in two weeks.
\end{textebox}
\begin{enumerate}
\item The communal work took place during the rainy season.
\item The inhabitants brought their own tools.
\item The session lasted two hours.
\item The council gave the volunteers gloves and water.
\item The next session will take place the following day.
\end{enumerate}
\end{exercice}

\begin{exercice}[Vocabulaire : trouvez le mot]
Complétez chaque phrase avec le mot anglais qui convient, choisi dans la liste : \eng{wheelbarrow} — \eng{gutters} — \eng{volunteer} — \eng{to provide} — \eng{neighbourhood} — \eng{tools}.
\begin{enumerate}
\item A \ldots\ldots\ldots\ldots is a person who works freely to help the community.
\item We use a \ldots\ldots\ldots\ldots to carry sand and stones.
\item After the rains, the \ldots\ldots\ldots\ldots were full of dirty water.
\item The council will \ldots\ldots\ldots\ldots food and drinks for the workers.
\item Everybody in our \ldots\ldots\ldots\ldots joined the clean-up campaign.
\item Hoes, cutlasses and brooms are useful \ldots\ldots\ldots\ldots for communal work.
\end{enumerate}
\end{exercice}

\begin{exercice}[Grammaire : complétez au bon temps]
Mettez les verbes entre parenthèses au \eng{simple past} ou au \eng{present perfect}, selon le sens.
\begin{enumerate}
\item Last Saturday, we \ldots\ldots\ldots\ldots (clean) the gutters of our street.
\item The volunteers \ldots\ldots\ldots\ldots (already / finish) the work.
\item The quarter head \ldots\ldots\ldots\ldots (announce) the date yesterday.
\item I \ldots\ldots\ldots\ldots (never / take part) in communal work before.
\item They \ldots\ldots\ldots\ldots (bring) their wheelbarrows two hours ago.
\end{enumerate}
\end{exercice}

\courssub{Je m'entraîne}

\begin{exercice}[Traduction]
Traduisez les phrases suivantes en anglais.
\begin{enumerate}
\item Les habitants ont nettoyé le marché samedi dernier.
\item Plus de cent volontaires ont participé aux travaux.
\item La mairie a fourni des gants et de l'eau potable.
\item Nous avons commencé à six heures du matin et fini à midi.
\item Le travail communautaire unit les gens du quartier.
\end{enumerate}
\end{exercice}

\begin{exercice}[Transformation de phrases]
Transformez les phrases selon la consigne entre parenthèses.
\begin{enumerate}
\item The volunteers cleaned the market. (mettre à la forme négative)
\item The council provided gloves for the workers. (mettre à la forme interrogative)
\item Eighty people took part in the work. (poser la question sur \eng{eighty} avec \eng{how many})
\item The quarter head said: “We will build a bridge.” (mettre au discours indirect)
\item They started the work at six o'clock. (poser la question avec \eng{when})
\end{enumerate}
\end{exercice}

\begin{exercice}[Texte à trous]
Complétez le texte suivant avec dix mots du glossaire de la leçon : \eng{volunteers} — \eng{wheelbarrow} — \eng{gutters} — \eng{provided} — \eng{communal work} — \eng{tools} — \eng{quarter head} — \eng{clean} — \eng{community} — \eng{lasted}.
\begin{textebox}[A clean-up day in Douala]
Last Sunday, the people of Bonabéri organised a \ldots\ldots\ldots\ldots session to \ldots\ldots\ldots\ldots the streets. About sixty \ldots\ldots\ldots\ldots came with their \ldots\ldots\ldots\ldots: hoes, brooms and one old \ldots\ldots\ldots\ldots to carry the rubbish. They cleared the \ldots\ldots\ldots\ldots which were blocked by plastic bags. The \ldots\ldots\ldots\ldots coordinated the teams, and the council \ldots\ldots\ldots\ldots water and biscuits. The work \ldots\ldots\ldots\ldots five hours. Everybody agreed that such activities make the \ldots\ldots\ldots\ldots stronger.
\end{textebox}
\end{exercice}

\begin{exercice}[Appariement : les outils]
Associez chaque outil anglais à sa traduction française, puis utilisez quatre de ces mots dans des phrases complètes de votre composition.
\begin{center}
\begin{tabularx}{\linewidth}{|X|X|}\hline
\rowcolor{popblueL} \textbf{English} & \textbf{Français} \\ \hline
1. a wheelbarrow & a) une pelle \\ \hline
2. a hoe & b) une machette \\ \hline
3. a cutlass & c) une brouette \\ \hline
4. a broom & d) une houe \\ \hline
5. a shovel & e) un râteau \\ \hline
6. a rake & f) un balai \\ \hline
7. gloves & g) un seau \\ \hline
8. a bucket & h) des gants \\ \hline
\end{tabularx}
\end{center}
\end{exercice}

\courssub{Je me prépare au Bac}

\begin{exercice}[Compréhension de l'écoute — type Bac]
Votre professeur vous lit (ou vous écoutez) le script suivant, \textbf{deux fois}. Prenez des notes, puis répondez aux questions.
\begin{textebox}[Building a community bridge in Nkambe]
Good morning, dear listeners. Today, our reporter takes us to Nkambe, in the North West Region, where the inhabitants have just finished building a wooden bridge over the small river that separates two quarters. The project lasted three weeks. Every Saturday, about one hundred and fifty volunteers — men, women and even students on holiday — gathered at seven o'clock in the morning. They carried stones, cut planks and fixed them with nails. The local carpenter, Mr Shey, directed the work for free. The council provided cement and iron bars, while a businessman from Bamenda offered ten bags of nails and two wheelbarrows. The quarter head, Mrs Wirba, told our reporter: “Before, children could not go to school when it rained. Now they cross the river safely every day.” The community plans to organise a big celebration next month to thank all the volunteers.
\end{textebox}
\textbf{A. Vrai ou faux justifié.} Dites si les affirmations sont \eng{true} ou \eng{false} et justifiez par une courte citation du script.
\begin{enumerate}
\item The bridge is made of iron.
\item The volunteers worked every day of the week.
\item Mr Shey was paid to direct the work.
\item A businessman offered nails and wheelbarrows.
\item Children can now go to school even when it rains.
\end{enumerate}
\textbf{B. Relevez les informations.} Notez les chiffres et noms propres entendus.
\begin{enumerate}
\item How long did the project last?
\item How many volunteers gathered every Saturday?
\item At what time did the work start?
\item Who is Mrs Wirba?
\item When will the celebration take place?
\end{enumerate}
\end{exercice}

\begin{exercice}[Traduction — type Bac]
Traduisez en français le passage suivant, extrait d'un article sur le travail communautaire.
\begin{textebox}[For translation]
Communal work is an old tradition in many Cameroonian villages. Once a month, the inhabitants leave their farms and shops to work together for the whole community. They clean the markets, repair the roads and clear the gutters before the rainy season. Nobody is paid, but everybody benefits: the village becomes cleaner, diseases disappear, and neighbours become friends. This spirit of solidarity is one of the greatest riches of our country.
\end{textebox}
\end{exercice}

\begin{exercice}[Expression écrite — type Bac]
\textbf{Sujet.} Your school wants to organise a clean-up day in your neighbourhood market. As the president of the English Club, write an announcement (about 80--100 words) inviting all the inhabitants to take part in the communal work. Your announcement must include the following five pieces of information:
\begin{enumerate}
\item what the activity is (cleaning the market);
\item the exact date and the starting time;
\item the meeting place;
\item the tools participants should bring;
\item what the organisers will provide (water, gloves, food).
\end{enumerate}
Use the vocabulary of the lesson (\eng{volunteers, communal work, gutters, wheelbarrow, provide}\ldots) and at least two sentences in the future with \eng{will}.
\end{exercice}



\newpage

\coursec{Corrigés des exercices}

\begin{corrige}[Exercice 1 — QCM de compréhension]
\begin{enumerate}
\item \textbf{b)} — \eng{“You are the quarter head of Ntarikon.”} Le journaliste s'adresse à M. Ndi en ces termes.
\item \textbf{c)} — \eng{“We cleaned the main market and cleared the gutters around it.”}
\item \textbf{c)} — \eng{“About eighty volunteers.”} Attention à l'écoute : \eng{eighty} (80) se prononce « èi-ti », alors que \eng{eighteen} (18) se prononce « èi-tine » avec l'accent sur la dernière syllabe.
\item \textbf{a)} — \eng{“We started at six o'clock in the morning.”} \eng{Noon} (midi) est l'heure de la fin, pas du début.
\item \textbf{b)} — \eng{“A clean market means fewer diseases, and it brings the community together.”}
\end{enumerate}
\textbf{Points de vigilance :} ne pas confondre les nombres en \emph{-ty} et en \emph{-teen} ; repérer les mots interrogatifs (\eng{who, what, how many, when, why}) pour cibler l'information attendue.
\end{corrige}

\begin{corrige}[Exercice 2 — Vrai ou faux justifié]
\begin{enumerate}
\item \textbf{False} — \eng{“to clean the streets after the rainy season”} : les travaux ont eu lieu \emph{après} la saison des pluies, pas pendant.
\item \textbf{True} — \eng{“inhabitants came with their own tools: hoes, cutlasses, wheelbarrows and brooms.”}
\item \textbf{False} — \eng{“The work lasted four hours, from seven to eleven in the morning.”} (quatre heures, pas deux)
\item \textbf{True} — \eng{“The council provided gloves and drinking water for everyone.”}
\item \textbf{False} — \eng{“the next session would take place in two weeks”} : dans deux semaines, pas le lendemain.
\end{enumerate}
\textbf{Points de vigilance :} les pièges portent sur les prépositions de temps (\eng{after} ≠ \eng{during}) et sur les durées ; toujours citer le texte pour justifier.
\end{corrige}

\begin{corrige}[Exercice 3 — Vocabulaire : trouvez le mot]
\begin{enumerate}
\item \eng{volunteer} (un volontaire, un bénévole)
\item \eng{wheelbarrow} (une brouette ; prononciation « ouil-bâ-ro »)
\item \eng{gutters} (les caniveaux ; pluriel imposé par le verbe \eng{were})
\item \eng{provide} (fournir ; après l'auxiliaire \eng{will}, le verbe reste à la base verbale, sans \emph{-s} ni \emph{-ed})
\item \eng{neighbourhood} (le quartier ; orthographe britannique avec \emph{-our-})
\item \eng{tools} (les outils ; pluriel imposé par \eng{are})
\end{enumerate}
\end{corrige}

\begin{corrige}[Exercice 4 — Grammaire : complétez au bon temps]
\begin{enumerate}
\item \eng{cleaned} — \eng{last Saturday} est un repère de temps passé révolu : \emph{simple past}.
\item \eng{have already finished} — \eng{already} appelle le \emph{present perfect} (action accomplie dont le résultat compte maintenant).
\item \eng{announced} — \eng{yesterday} impose le \emph{simple past}.
\item \eng{have never taken part} — \eng{never\ldots before} exprime une expérience de vie : \emph{present perfect} ; participe passé irrégulier \eng{take -- took -- taken}.
\item \eng{brought} — \eng{two hours ago} impose le \emph{simple past} ; verbe irrégulier \eng{bring -- brought -- brought}.
\end{enumerate}
\textbf{Points de vigilance :} en anglais, on n'utilise \emph{jamais} le \emph{present perfect} avec un repère de temps passé précis (\eng{yesterday, last week, ago}) — contrairement au français « j'ai fait hier ».
\end{corrige}

\begin{corrige}[Exercice 5 — Traduction]
\begin{enumerate}
\item \eng{The inhabitants cleaned the market last Saturday.} (repère passé : \emph{simple past})
\item \eng{More than one hundred volunteers took part in the work.} (\eng{to take part in} = participer à ; la préposition \eng{in} est obligatoire)
\item \eng{The council provided gloves and drinking water.} (\eng{the council} = la mairie, le conseil municipal ; \eng{drinking water} = eau potable)
\item \eng{We started at six o'clock in the morning and finished at noon.}
\item \eng{Communal work unites the people of the neighbourhood.} (vérité générale : \emph{present simple}, avec \emph{-s} à la 3\textsuperscript{e} personne du singulier)
\end{enumerate}
\textbf{Points de vigilance :} \eng{people} est un pluriel (« les gens ») : \eng{the people unite}, mais ici le sujet est \eng{communal work} (singulier), d'où \eng{unites}.
\end{corrige}

\begin{corrige}[Exercice 6 — Transformation de phrases]
\begin{enumerate}
\item \eng{The volunteers did not clean the market.} (négation au \emph{simple past} : \eng{did not} + base verbale ; le verbe principal perd sa marque \emph{-ed})
\item \eng{Did the council provide gloves for the workers?} (interrogation : \eng{did} + sujet + base verbale)
\item \eng{How many people took part in the work?} (\eng{how many} + nom pluriel)
\item \eng{The quarter head said (that) they would build a bridge.} (discours indirect : \eng{will} devient \eng{would} ; \eng{we} devient \eng{they})
\item \eng{When did they start the work?}
\end{enumerate}
\textbf{Points de vigilance :} après \eng{did} ou \eng{did not}, toujours la base verbale — jamais \eng{did not cleaned}.
\end{corrige}

\begin{corrige}[Exercice 7 — Texte à trous]
Dans l'ordre des blancs : \eng{communal work} — \eng{clean} — \eng{volunteers} — \eng{tools} — \eng{wheelbarrow} — \eng{gutters} — \eng{quarter head} — \eng{provided} — \eng{lasted} — \eng{community}.\\
Texte complété : \eng{“Last Sunday, the people of Bonabéri organised a communal work session to clean the streets. About sixty volunteers came with their tools: hoes, brooms and one old wheelbarrow to carry the rubbish. They cleared the gutters which were blocked by plastic bags. The quarter head coordinated the teams, and the council provided water and biscuits. The work lasted five hours. Everybody agreed that such activities make the community stronger.”}
\end{corrige}

\begin{corrige}[Exercice 8 — Appariement : les outils]
1--c ; 2--d ; 3--b ; 4--f ; 5--a ; 6--e ; 7--h ; 8--g.\\
Exemples de phrases :
\begin{itemize}
\item \eng{The volunteers carried the sand in a wheelbarrow.} « Les volontaires ont transporté le sable dans une brouette. »
\item \eng{She swept the street with a broom.} « Elle a balayé la rue avec un balai. »
\item \eng{We cut the grass with cutlasses.} « Nous avons coupé l'herbe avec des machettes. »
\item \eng{The council gave gloves to all the workers.} « La mairie a donné des gants à tous les travailleurs. »
\end{itemize}
\textbf{Points de vigilance :} \eng{a cutlass} = une machette (faux ami : ce n'est pas « un couteau ») ; \eng{a hoe} se prononce « ro » (le \emph{h} est muet).
\end{corrige}

\begin{corrige}[Exercice 9 — Compréhension de l'écoute — type Bac]
\textbf{A. Vrai ou faux justifié.}
\begin{enumerate}
\item \textbf{False} — \eng{“building a wooden bridge”} : le pont est en bois, pas en fer.
\item \textbf{False} — \eng{“Every Saturday\ldots gathered at seven o'clock”} : les volontaires se réunissaient chaque samedi, pas tous les jours.
\item \textbf{False} — \eng{“directed the work for free”} : gratuitement, donc sans être payé.
\item \textbf{True} — \eng{“a businessman from Bamenda offered ten bags of nails and two wheelbarrows.”}
\item \textbf{True} — \eng{“Now they cross the river safely every day.”}
\end{enumerate}
\textbf{B. Relevez les informations.}
\begin{enumerate}
\item \eng{The project lasted three weeks.}
\item \eng{About one hundred and fifty volunteers gathered every Saturday.}
\item \eng{The work started at seven o'clock in the morning.}
\item \eng{Mrs Wirba is the quarter head.}
\item \eng{The celebration will take place next month.}
\end{enumerate}
\textbf{Barème indicatif (10 pts) :} partie A : 5 × 1 pt (0,5 pt pour la réponse + 0,5 pt pour la justification) ; partie B : 5 × 1 pt.\\
\textbf{Points de vigilance :} à l'écoute, noter d'abord les chiffres (150, 7 h, 3 semaines) et les noms propres (Nkambe, M. Shey, Mme Wirba, Bamenda) ; \eng{for free} = « gratuitement », piège classique.
\end{corrige}

\begin{corrige}[Exercice 10 — Traduction — type Bac]
Proposition de traduction :\\
« Le travail communautaire est une vieille tradition dans de nombreux villages camerounais. Une fois par mois, les habitants quittent leurs champs et leurs boutiques pour travailler ensemble pour toute la communauté. Ils nettoient les marchés, réparent les routes et débouchent les caniveaux avant la saison des pluies. Personne n'est payé, mais tout le monde en profite : le village devient plus propre, les maladies disparaissent et les voisins deviennent des amis. Cet esprit de solidarité est l'une des plus grandes richesses de notre pays. »\\
\textbf{Barème indicatif (5 pts) :} 1 pt par segment de sens correctement rendu (5 segments).\\
\textbf{Points de vigilance :} \eng{once a month} = « une fois par mois » ; \eng{benefits} est ici un \emph{verbe} (« profite »), pas le nom « bénéfices » ; \eng{farms} = « champs/exploitations », pas « fermes » au sens bâtiment ; \eng{neighbours} s'écrit avec \emph{-our-} en anglais britannique.
\end{corrige}

\begin{corrige}[Exercice 11 — Expression écrite — type Bac]
Modèle de rédaction (environ 95 mots) :
\begin{quote}
\eng{\textbf{ANNOUNCEMENT — CLEAN-UP DAY!}\\
Dear inhabitants of Mokolo,\\
The English Club of our school is organising a communal work session to clean Mokolo market on Saturday, 15th June, starting at 6.30 a.m. We shall meet in front of the main entrance of the market. Please bring your own tools: brooms, hoes, cutlasses, buckets and wheelbarrows. We will clear the gutters and sweep the whole market. The organisers will provide gloves, drinking water and food for all the volunteers. Come in large numbers! Together, we will make our neighbourhood cleaner and healthier.\\
The President of the English Club}
\end{quote}
\textbf{Barème indicatif (10 pts) :} contenu (5 informations demandées) : 5 pts ; réemploi du lexique de la leçon : 2 pts ; correction grammaticale (dont deux phrases au futur avec \eng{will}) : 2 pts ; présentation de l'annonce (titre, date, signature) : 1 pt.\\
\textbf{Points de vigilance :} date à l'anglaise (\eng{Saturday, 15th June}) ; \eng{at} devant l'heure, \eng{on} devant le jour ; ne pas oublier le lieu de rendez-vous (\eng{meeting place}).
\end{corrige}

\newpage

\coursec{Fiche bilan}

\begin{popfigure}
\begin{tikzpicture}[mindmap, grow cyclic, every node/.style={concept}, text=white,
root concept/.append style={concept color=popblue, font=\bfseries},
level 1/.append style={level distance=3.4cm, sibling angle=51.4},
level 2/.append style={level distance=2.1cm}]
\node[root concept, align=center]{Listening\\ Communal\\ Work}
child[concept color=poporange]{ node[align=center]{Theme\\ community\\ work}
  child { node[concept color=poporange!70, align=center]{clean-up\\ campaign} }
  child { node[concept color=poporange!70, align=center]{village\\ meeting} }
}
child[concept color=popgreen]{ node[align=center]{Keywords\\ to catch}
  child { node[concept color=popgreen!70, align=center]{verbs:\\ join, help, dig} }
}
child[concept color=poppurple]{ node[align=center]{Numbers\\ and names}
  child { node[concept color=poppurple!70, align=center]{dates, times,\\ places} }
}
child[concept color=poppink]{ node[align=center]{Grammar\\ present simple\\ past simple}
  child { node[concept color=poppink!70, align=center]{frequency\\ adverbs} }
}
child[concept color=popgold]{ node {Question types}
  child { node[concept color=popgold!70] {true / false} }
  child { node[concept color=popgold!70, align=center]{MCQ, short\\ answers} }
}
child[concept color=popblue!80]{ node {Strategy}
  child { node[concept color=popblue!60, align=center]{read questions\\ first} }
  child { node[concept color=popblue!60] {listen twice} }
};
\end{tikzpicture}
\legende{Carte mentale de la leçon : écouter un texte sur le travail communautaire}
\end{popfigure}

\begin{aretenir}[L'essentiel de la leçon — 1. Lexique clé du travail communautaire]
\begin{itemize}
\item \cle{communal work} : travail communautaire ; \cle{community} : communauté
\item \cle{to take part in} : participer à ; \cle{to join} : rejoindre, se joindre à
\item \cle{clean-up campaign} : campagne de nettoyage ; \cle{to sweep} : balayer
\item \cle{to dig} : creuser ; \cle{to plant trees} : planter des arbres
\item \cle{to collect rubbish} : ramasser les ordures ; \cle{rubbish} : ordures (nom indénombrable)
\item \cle{neighbourhood} : quartier ; \cle{village chief} : chef de village
\item \cle{volunteer} : bénévole (nom) ; \cle{to volunteer} : se porter volontaire (verbe)
\item \cle{to cooperate} : coopérer ; \cle{teamwork} : travail d'équipe
\item \cle{wheelbarrow} : brouette ; \cle{hoe} : houe ; \cle{machete} : machette
\item \cle{to repair a bridge} : réparer un pont ; \cle{well} : puits
\end{itemize}
\end{aretenir}

\begin{propriete}[L'essentiel de la leçon — 2. Grammaire à connaître par c\oe ur]
\begin{center}
\begin{tabularx}{\linewidth}{|X|X|X|}
\hline
\rowcolor{popblueL} Point & Forme & Exemple \\
\hline
Present simple (habitudes) & sujet + verbe (3\textsuperscript{e} pers. du singulier : -s) & \eng{The village meets every Saturday.} « Le village se réunit chaque samedi. » \\
\hline
Past simple (faits passés) & sujet + verbe en -ed / 2\textsuperscript{e} colonne des verbes irréguliers & \eng{They cleaned the market last week.} « Ils ont nettoyé le marché la semaine dernière. » \\
\hline
Adverbes de fréquence & avant le verbe principal, après le verbe \eng{be} & \eng{We always help our neighbours.} « Nous aidons toujours nos voisins. » \\
\hline
\eng{take part in} + nom & jamais \eng{take part to} & \eng{She took part in the clean-up.} « Elle a participé au nettoyage. » \\
\hline
\end{tabularx}
\end{center}
\end{propriete}

\begin{attention}[Verbes irréguliers fréquents du thème — à vérifier à l'écrit]
\begin{center}
\begin{tabular}{|l|l|l|}
\hline
\rowcolor{popblueL} Base verbale & Past simple & Traduction \\
\hline
take & took & prendre, participer (\eng{take part in}) \\
\hline
sweep & swept & balayer \\
\hline
dig & dug & creuser \\
\hline
bring & brought & apporter \\
\hline
build & built & construire \\
\hline
meet & met & se réunir, rencontrer \\
\hline
\end{tabular}
\end{center}
Attention : \eng{sweep} ne fait pas \emph{sweeped} mais \eng{swept} ; \eng{dig} ne fait pas \emph{digged} mais \eng{dug}.
\end{attention}

\begin{methode}[L'essentiel de la leçon — 3. Méthode express d'écoute]
\begin{enumerate}
\item \textbf{Avant l'écoute} : lis les questions, souligne les mots-clés (noms propres, chiffres, verbes d'action), prédis le thème.
\item \textbf{Première écoute} : note le thème général, les noms propres, les chiffres, les dates et les lieux.
\item \textbf{Deuxième écoute} : vérifie les détails, complète les réponses courtes, tranche les vrai/faux.
\item \textbf{Après l'écoute} : relis tes réponses (accord sujet--verbe, orthographe des verbes irréguliers, logique de la réponse).
\end{enumerate}
\end{methode}

\begin{exemplebox}[L'essentiel de la leçon — 4. Expressions utiles pour réemployer le lexique]
\begin{itemize}
\item \eng{Last Saturday, our community took part in a clean-up campaign.} « Samedi dernier, notre communauté a participé à une campagne de nettoyage. »
\item \eng{Many volunteers swept the streets and collected rubbish.} « De nombreux bénévoles ont balayé les rues et ramassé les ordures. »
\item \eng{The village chief asked everyone to bring a hoe or a machete.} « Le chef de village a demandé à chacun d'apporter une houe ou une machette. »
\item \eng{We usually meet at the well at seven o'clock in the morning.} « Nous nous retrouvons habituellement au puits à sept heures du matin. »
\end{itemize}
\end{exemplebox}

\begin{aretenir}[Checklist avant le Bac]
\begin{itemize}
\item[$\square$] Je connais le lexique du travail communautaire (au moins 15 mots).
\item[$\square$] Je sais repérer le thème général d'un enregistrement dès la première écoute.
\item[$\square$] Je lis les questions avant d'écouter et je souligne les mots-clés.
\item[$\square$] Je relève les chiffres, dates, heures et noms propres sans erreur.
\item[$\square$] Je distingue present simple (habitude) et past simple (fait passé).
\item[$\square$] Je place correctement les adverbes de fréquence (\eng{always, usually, sometimes}).
\item[$\square$] Je dis \eng{take part in} et jamais \eng{take part to}.
\item[$\square$] Je sais répondre aux vrai/faux en justifiant par une phrase du texte.
\item[$\square$] Je réponds aux QCM en éliminant les distracteurs.
\item[$\square$] Je rédige des réponses courtes grammaticalement correctes.
\item[$\square$] Je réemploie le lexique dans trois phrases personnelles.
\item[$\square$] Je vérifie l'orthographe des verbes irréguliers au passé (\eng{took, swept, dug}).
\end{itemize}
\end{aretenir}

\findecours

\end{document}
